% 本文件由 LaTeX 排版 Agent 根据有效 translation.json 自动生成。
% author=Ambrose; work_id=3407; title=论童贞

\chapter{论童贞}

% structure: p0001 source=h1 target=omitted reason=page-title-already-rendered-by-parent

卷一\newline{}卷二\newline{}卷三\par

\SourceRecord{Translated by H. de Romestin, E. de Romestin and H.T.F. Duckworth.；Nicene and Post-Nicene Fathers, Second Series；Vol. 10.；Edited by Philip Schaff and Henry Wace.；Buffalo, NY: Christian Literature Publishing Co.,；1896.；<http://www.newadvent.org/fathers/3407.htm>.}

% structure: p0001 source=h1 target=section reason=page-title

\section{论童贞（卷一）}

% structure: p0002 source=h2 target=subsection reason=relative-heading-rank; base=h2

\subsection{引言}

\OriginalTerm{童贞}{Virginity}的状态无疑在\OriginalTerm{圣经}{holy Scripture}中受到称赞，既由主也由\Person{圣保禄}；但学者们对最初时代贞女们所遵循的习俗与规矩，意见分歧。有些人认为，从一开始，贞女们便有习惯公开\OriginalTerm{宣发}{profession}贞洁圣愿，并共同生活。另一些人则认为，她们的圣愿是私下的，她们有时共同生活，有时住在父母家里。还有一些人相信，贞女们的意向仅以面纱和衣着朴素来表示，而团体生活的最早开始应归功于\Person{圣盎博罗削}本人。\par

第一种意见，就任何为人所知的宣发圣愿而言，是难以成立的。早期\WorkTitle{殉道者行传}中的记载应持怀疑态度，因为这类著作中多有伪造。教父们和公会议的言论，在这一点上，所能确立的也几乎不超过上述第二种意见。\par

似乎曾有一些人，如\Person{圣盎博罗削}的姐妹\Person{玛尔切利纳}，公开宣发了圣愿，并由主教手中正式领受了面纱；另有同样志向坚定的人，她们的童贞圣愿是私下发出的。前者中，住在\Place{米兰}的那些人似乎几乎未过团体生活，但在\Place{博洛尼亚}[I. 60]她们却如此。\OriginalTerm{圣愿}{vow}、\OriginalTerm{领受面纱}{taking the veil}和\OriginalTerm{宣发圣愿}{profession}这些术语，在\Person{圣盎博罗削}的时代与现在一样使用。\par

那么，看来从宗徒时代起，就有人以贞洁的生活将自己奉献给\OriginalTerm{天主}{God}，后来这一许诺或圣愿是在他人——主教、圣职人员和朋友面前发出的。迫害延续的时期，这些贞女与父母住在家里，因为她们实际上不可能住到别处。他们之间的团体生活似乎始于\Place{东方}，而当\Person{圣亚大纳修斯}为躲避亚略派而来到\Place{罗马}时，便将这一习俗引入西方教会。\par

\Person{圣盎博罗削}沿着这一方向积极工作，不仅在自己的教区，还在邻近的省份，甚至在\Place{非洲}。在主教任期之初，他就此主题向自己的羊群讲话，并应他姐妹\Person{玛尔切利纳}的请求，将他的教导汇编成了以下三卷书。\par

在第一卷，他论及\OriginalTerm{童贞}{Virginity}的尊严，并说明他写作的原因。当他在\Person{圣依搦斯}殉道纪念日开始其讲话时，便以她的故事作为该论著前部的主题，并说明在犹太人中间、甚至在外教人中间，童贞的恩宠是如何被预示，并最终因我主的来临而被宣告的。然后他警告父母，尤其是寡妇，不要阻止他们的女儿聆听关于这一主题的讲话，并提及那些甚至从很远的地方来到\Place{米兰}领受面纱的人数。\par

在第二卷，在论及贞女的品格与生活方式时，他正如自己所说，是通过列举榜样和实例，而非制定一套规则来进行的。他谈到\Person{德克拉}，\Place{米兰}的主保圣人，\Person{圣保禄}的门徒，以及其他贞女。\par

在第三卷，他回顾了\Person{教宗利伯略}在一次大会众前亲手授予\Person{玛尔切利纳}面纱时所作讲话的概要。\Person{圣盎博罗削}提出了一些告诫，反对过分的刻苦，并建议以祈祷和修习内心德行，取代某些外在行为。针对\Person{玛尔切利纳}的一个问题，他还详细讨论了某些宁可自杀也不愿失贞的贞女这一话题。\par

作者本人说，这部论著写成于其主教任期第三年，即公元377年，并得到\Person{圣热罗尼莫}的赞同引用，见\WorkTitle{书信集}XXII.22和XLVIII.14[本系列卷六，第31及75页]，以及\Person{圣奥斯定}，\WorkTitle{论基督教导}IV.48，50。\par

\Emph{致他的姐妹玛尔切利纳}\par

% structure: p0013 source=h2 target=subsection reason=relative-heading-rank; base=h2

\subsection{第一章}

\Emph{\Person{圣盎博罗削}，思考了他必须就自己所领受的塔冷通交账这事，便决定写作，并以天主慈悲的某些例证安慰自己。然后，认识到自己的不足后，他希望能像福音中的无花果树那样被对待，并表达希望在努力宣讲基督时不致词穷。}\par

1. 如果按照天上真理的定断，我们必须要为说过的每一句闲话交账，\BibleRef{玛 12:36}，如果每个仆人，当他的主人回来时，若因他像胆小的放债人或贪婪的物主，将委托给他、本应通过利息而增值的\OriginalTerm{属神的恩宠}{spiritual grace}的\OriginalTerm{塔冷通}{talents}埋藏在地里，便要受到不小的责备，那么我——虽只有中等才能，却有极大的责任使托付给我的天主圣言增值——自然应该担心，生怕要我来交代我话语的益处，尤其因为主向我们索要的是努力，而非成果。因此，我决定写下些什么，更何况我的话语被人聆听比写下时更冒有失谦逊的风险，因为书不知谦逊。\par

2. 因此，我虽实不信任己能，却因\OriginalTerm{天主慈悲}{divine mercy}的实例而受鼓舞，大胆草拟一篇讲词，因为当天主愿意时，连驴也开口说了话。\BibleRef{户 22:28} 我将开启久已缄默的口，愿天使也来扶助我，使我摆脱此世的牵累，因为那位在驴身上除去了本性障碍者，也能除去我拙于言辞的障碍。在旧约的\OriginalTerm{约柜}{ark}中，司祭的手杖发了芽；\BibleRef{户 17:8} 在天主，在\OriginalTerm{圣教会}{Holy Church}内从我们的结节处也生出一朵花来，是何等容易。我们又怎能怀疑，那位曾在荆棘丛中说话的天主，不会在人内说话呢？\BibleRef{出 3:4} 天主并未轻视荆棘，但愿他也赐予我的荆棘以光明。或许有人会惊奇，在我们的荆棘中竟也有光明；有些荆棘不会被烧毁；有些人一听到我的声音就会脱去脚上的鞋，好让心灵的步履摆脱肉身的羁绊。\par

3. 但是这些事乃是圣善的人所获得的。愿耶稣垂顾仍躺在那棵不结果的无花果树下的我\BibleRef{若 1:48}，并且我的无花果树三年后也能结果实。然而罪人怎能有如此大的盼望呢？但愿至少那福音中葡萄园的园丁，或许已受命砍掉我的无花果树，今年仍容它留着，直到他周围掘土、加上粪肥，或许他能从尘土中扶起无助者，从粪土中举起贫穷者。那些把马拴在葡萄树和橄榄树下的人是有福的\BibleRef{创 49:11}，他们将自己的辛劳之旅奉献给光明与喜乐；而无花果树，即世俗享乐的诱惑，仍然笼罩着我，低矮易折，柔软无用，不结果实。\par

4. 有人或许诧异，我既不能言，竟敢写作。然而如果我们思想福音书中记载的，以及司铎们的行事，并以圣先知匝加利亚为例，他就会发现有些事情是声音无法解释的，但笔却能写出来。如果若翰的名字使他父亲恢复了说话的能力\BibleRef{路 1:63}，那么我也不可绝望：我虽哑，但只要讲论基督，仍可获得言语；论到基督，正如先知的话所说：\Quote{谁将述说祂的世代？}\BibleRef{依 53:8} 所以我作为仆人，要宣布主的家族，因为主在这个充满软弱的人性肉躯内，为自己祝圣了一个家族。\par

% structure: p0019 source=h2 target=subsection reason=relative-heading-rank; base=h2

\subsection{第二章。}

\Emph{本论著有一个良好的开始，因为今天是圣童贞依搦斯的生日，圣安博论及她的名字、端庄和殉道，加以赞扬，但尤其论及她的年龄，看她虽然年仅十二岁，却能以全副男儿气概，胜过恐吓、许诺、酷刑和死亡本身。}\par

5. 我的任务开始得良好，因为今天是一位童贞的生日，我须讲论童贞，而这论著便由此篇讲道开始。今天是殉道者的生日，让我们献上祭品。今天是圣依搦斯的生日，让男子们钦佩，让孩子们鼓起勇气，让已婚者惊奇，让未婚者效法。但我能说什么配得上她呢？她的名字本身就非乏善可陈。在超过年龄的热忱中，在超性的德行上，她在我看来，似乎担当的不是一个人间的名字，而是一个殉道的表记，借此她显明了她将成为怎样的人。\par

6. 然而我有可助我者。童贞之名，便是端庄的称号。我将呼求殉道者，宣扬童贞。那篇颂词无需铺陈，伸手可及，已够冗长。让劳苦停止吧，让雄辩缄默吧。一句话便足以赞美。这句话老幼青壮都在吟唱。无人比那受众人赞美者更值得称扬。有多少人，就有多少传报者，当他们开口时，便是宣布殉道者。\par

7. 据说她受难殉道时年方十二。那残忍更加可憎，因竟不放过如此稚弱之龄，而信德的力量实则更显伟大，即使在此年龄中亦获得了证据。在那纤小的身躯上，容得下伤口吗？而那甚至无处置剑锋之躯，却有法子胜过剑锋。然而这般年岁的少女，连父母的怒容亦承受不住，常因针刺般的微小痛楚而哭泣，如同受伤一般。在残酷的刽子手手下，她毫无畏惧；在沉重作响的铁链重压下，她不为所动，将自己整个身躯献予狂暴士兵的刀剑，虽尚不知死亡为何，却已预备好了。若她被强行拖至祭台，她也预备好在祭火前向基督伸出双手，在亵圣的祭台前划上主、征服者的十字圣号，或又将颈项和双手伸入铁枷，但无枷锁能困住如此纤细的肢体。\par

8. 新式的殉道！尚未达受刑之龄，却已成熟以待凯旋，难以争战却易得荣冠，她虽身处年幼之不利，却履行了教授勇德之职。她不似新娘奔赴婚床，身为童贞女，反急步欣然往受苦之地，头上所饰的不是编结的发辫，而是基督。众人都流泪，唯她无泪。众人惊奇她如此轻易地挥霍尚未享用的生命，现在却放弃，仿佛已经历过一般。人人都诧异，竟有人能为天主作证，而此人因年龄尚幼，尚不能自主。她竟使人相信她关于天主的证言，而关于人的证据却无人接受。因为超越本性之事，乃出自本性之造主。\par

9. 刽子手用何等的威胁使她惧怕，用何等的诱惑劝诱她，多少人想要娶她为妻！但她回答说：\Quote{如果我注视任何我以为能取悦于我的人，就是对我的净配的侮辱。那位首先为自己拣选我者，将接纳我。刽子手，你为何迟延？让这躯体消亡吧，它可能被我不愿见的眼睛所喜爱。}她站着，祈祷，俯首。你可见刽子手战抖，仿佛他自己被定了罪，右手发颤，面色苍白，他惧怕别人的危难，而少女却不为自己惧怕。因此，在一个牺牲身上，你有了双重的殉道：端庄的殉道和信仰的殉道。她既保持了童贞，又获得了殉道。\par

% structure: p0026 source=h2 target=subsection reason=relative-heading-rank; base=h2

\subsection{第三章。}

\Emph{童贞因多种理由受赞扬，但主要是因为它将圣言从天上带下，因此，对它的追求，在旧盟约下只有少数人实行，现已遍及无数人。}\par

10. 如今，对纯洁的爱催迫着我，而你，我的圣善姊妹，虽沉默的会衣未发一言，让我谈一谈童贞，免得这主要的德行似乎只被略略带过。因为童贞之值得赞扬，并非因它在殉道者身上找到，而是因为它本身造就殉道者。\par

但谁以人的理智理解那连自然都未纳入其法则的事呢？或者谁以日常语言解释那超越自然过程的事呢？童贞从天上带来了它在地上效法的事物。而且她并非不恰当地从天寻求生活方式，因为她在天上为自己找到了一位净配。她穿越云层、空气、天使和星辰，在圣父的怀抱中找到了圣言，并以全心吸引到自己内。因为谁找到如此伟大的善之后会舍弃它？因为\Quote{祢的名如同倒出来的香膏，所以众童女都爱祢，吸引祢。} 而我所言并非出于自己，因为那不婚不嫁的人就像天上的天使。那么，如果她们被比拟为天使，并联合于天使之主，我们不必惊讶。那么，谁能否认这种生活方式源于天上，我们在地上难以找到，除非因天主降临到属地肉身的肢体中？那时，一位童贞女受孕，圣言成了血肉，为使血肉能成为天主。\par

但有人会说：\Quote{但厄里亚似乎与肉体之爱的拥抱毫无关系。} 因此他被马车带到天上，因此他荣耀地与主一同显现，\BibleRef{玛 17:3}，因此他将作为主降临的前驱而来。\BibleRef{拉 4:5} 而米黎盎拿起铃鼓，以贞女的端庄带领舞蹈。\BibleRef{出 15:20} 但请思考她当时所预表的是谁。她岂不是教会的类型吗？教会如同贞女，以无玷的精神，联合民众的宗教集会，咏唱神圣的歌曲。因为我们读到在耶路撒冷的圣殿中也曾选派贞女。但宗徒说了什么？\BibleQuote{这些事发生在他们身上作为预象，为作将来之事的标识。} \BibleRef{格前 10:11} 因为预象显示在少数人身上，而生命存在于许多人中。\par

但实际上，自从主带着我们的肉身降来，以无混淆、无混合的方式结合了天主性与肉身之后，那天上生活的方式便遍布全世界，被植入人的身体。这就是那在地上服侍的天使所预示将要发生的事，\BibleRef{玛 4:11}，此种服侍应以无玷身体的侍奉献给主。这就是那欢跃的天使对大地所说的天上事奉。\BibleRef{路 2:13} 因此，我们有自古以来的古老权威，以及从基督本身而来的完全阐明。\par

% structure: p0032 source=h2 target=subsection reason=relative-heading-rank; base=h2

\subsection{第四章}

\Emph{童贞的美丽从未存在于异教徒中，无论是\OriginalTerm{维斯塔贞女}{Vestal Virgins}，还是哲学家，如\Person{毕达哥拉斯}。}\par

我\Emph{确实}与异教徒毫无共同之处，就此事而言，我既不同流于蛮族，也不与野兽为伍；我们虽与它们呼吸同一种生命气息，拥有共同的属地身体条件，且生育方式无异，但唯独在这一点上，我们避免了类似的非难：童贞为异教徒所追求，但一经祝圣便被亵渎；为蛮族所侵袭，而在其他族类中则一无所知。\par

谁会对我提起\OriginalTerm{维斯塔}{Vesta}的贞女和\OriginalTerm{帕拉斯}{Pallas}的祭司呢？那算什么贞洁，不是出于道德，而是按年龄，不是设定为永远，而是设定期限！那种纯洁更为放荡，其败坏只是推迟到晚年而已。他们教导贞女不应坚守，且无法坚守，因为他们为童贞设定了期限。那算什么宗教，要求端庄的少女成为不端庄的老妇？受法律约束的人并不端庄，而被法律释放的人则是不端庄的。啊，奥秘！啊，道德！那里法律强迫贞洁，却又为情欲授权！所以，受恐惧约束的不是贞洁者；以价格雇来的不是可敬者；那种暴露于淫邪目光的每日纠缠，被可耻注视攻击的，也算不上端庄。她们获得豁免，得到酬金，仿佛出售贞洁不是放荡的最大证据。为酬金承诺的，为酬金放弃；为酬金转让；被认为有其价格。习惯出售贞洁的，不知道如何赎回它。\par

对于\OriginalTerm{弗鲁吉亚}{Phrygian}的仪式，我能说什么呢？那里淫乱是准则，且是较弱的性别的准则。对于\OriginalTerm{巴克斯}{Bacchus}的狂欢，我还能说什么呢？那里仪式的奥秘是情欲的刺激。那么，在这些地方，祭司们的生活能是什么样子的？那里神明们的奸淫竟是宗教事务。所以他们根本没有神圣的贞女。\par

让我们看看是否哲学家的训诫或许塑造了一些贞女，因为他们常声称教授一切美德。故事中提到一位\OriginalTerm{毕达哥拉斯派}{Pythagorean}的贞女，一个暴君试图强迫她透露秘密，以免在酷刑中仍能被迫泄露，她咬下舌头，吐在暴君脸上，使那不肯停止审问的人再无东西可问。\par

但那位心灵如此坚定的贞女，却被情欲征服，尽管她不能被酷刑征服。因此，能保守心灵秘密的她，却不能隐藏身体的羞耻。她战胜了自然，却未遵守纪律。她多么希望自己的言语曾作为贞洁的防卫而存在！所以，她并非在各方面都不可战胜，因为虽然暴君未能发现他所寻找的，却发现了未曾寻找的。\par

19. 我们的贞女们是何等坚强啊！她们甚至战胜了那些看不见的势力；她们的胜利不仅针对血肉，更是针对这世界的首领和今世的主宰！就年龄而言，\Person{依搦斯}确实较小，但论德行却更伟大，她战胜了更多敌人，信心也更坚定；她没有因恐惧而毁掉自己的舌头，反而将它留作战利品。因为她身上没有什么害怕泄露的，因为她所承认的是圣洁的，而非罪恶的。因此，前者只是隐藏了自己的秘密，后者却为主作证，并以自己的身体宣认了祂——尽管她的年龄还不容许她宣认。\par

% structure: p0040 source=h2 target=subsection reason=relative-heading-rank; base=h2

\subsection{第五章}

\Emph{天堂是童贞的居所，天主圣子是其创始者；祂虽然在圣母之前已是童贞者，却由圣母取了童贞教会作自己的净配。我们都是由她所生。她的恩赐部分已列举。她的女儿们享有特殊的卓越之处，因为童贞并非来自诫命，且它是追求虔诚的最有力帮助。}\par

20. 颂词的习惯是谈论对象的祖国和出身，通过提及父亲来彰显后裔的伟大。如今，我并非要赞扬童贞，而是要阐明它，但仍认为说明其祖国和父亲是有意义的。首先，让我们确定它的祖国在哪里。如果说祖国就是出生之地，那么毫无疑问，贞洁的本乡就是天堂。因此，她在此地是异乡人，在那里却是居民。\par

21. 童贞的贞洁无非是无玷的纯洁。除了无玷的天主圣子，我们还能认为谁是它的创始者呢？祂的肉身未见腐朽，祂的天主性未受玷染。那么，请想想童贞的功绩何等伟大。基督先于圣母，基督亦由圣母所生。祂于万世之前确实由圣父所生，但为了万世却由圣母诞生。前者出于祂的本性，后者为了我们的益处。前者恒常存在，后者则是祂所意愿的。\par

22. 再想想童贞的另一功绩。基督是童贞的净配，如果我们可以这样论及童贞的贞洁，因为童贞属于基督，而非基督属于童贞。因此，祂是那位曾许配的童贞女，是那位生了我们的童贞女，以她自己的乳汁喂养我们；关于她我们读到：\Quote{耶路撒冷的童贞女成就了何等大事！岩石的乳头不会枯竭，黎巴嫩的雪不会断绝，被强风吹送的水也不会缺乏。}这位被天主圣三的溪流浇灌的童贞女是谁？从她的岩石流出水，她的乳头不枯竭，她的蜜倾流而出。如今，根据宗徒的话，岩石就是基督。\BibleRef{格前 10:4}因此，从基督那里，乳头不枯竭，从天主那里，光辉不暗淡，从圣神那里，河流不干涸。这就是浇灌他们教会的天主圣三：圣父、基督和圣神。\par

23. 现在让我们从母亲转向女儿们。\BibleQuote{论到贞女们，}宗徒说，\BibleQuote{我没有主的命令。}\BibleRef{格前 7:25}如果外邦人的老师都没有，谁还能有呢？事实上，他没有命令，但有榜样。因为童贞不能被命令，而必须被渴求，因为那些超越我们的事，更适合祈求，而非受控制。他又说：\BibleQuote{我愿你们无所挂虑。没有妻子的人，挂虑的是主的事，想怎样悦乐天主……童贞女挂虑的是主的事，为能在身体和灵魂上都是圣的；而已婚的妇女，挂虑的是世俗的事，想怎样悦乐丈夫。}\par

% structure: p0046 source=h2 target=subsection reason=relative-heading-rank; base=h2

\subsection{第六章}

\Emph{圣盎博罗削表明他并非反对婚姻，进而比较独身与已婚状态的利弊。}\par

24. 我确实并未劝阻婚姻，而是在阐扬童贞的益处。宗徒说：\BibleQuote{软弱的人，只吃蔬菜。}\BibleRef{罗 14:2}我认为有些事是必需的，我赞赏另一些事。\BibleQuote{你与妻子有了约束吗？不要寻求解脱。你脱离妻子了吗？不要寻求妻子。}\BibleRef{格前 7:27}这是给已有之人的命令。但他对贞女们怎么说呢？\BibleQuote{谁把自己的童贞女出嫁，做得好；谁不出嫁，做得更好。}\BibleRef{格前 7:38}一个若是结婚，并不犯罪；另一个，若不结婚，却是为了永恒。前者是软弱的补救，后者是贞洁的光荣。前者不受责备，后者却受称赞。\par

25. 如果你们愿意，我们来比较一下已婚妇女的益处与等待贞女的福分。虽然高贵的妇女夸耀自己子女众多，但她生得越多，忍受得也越多。让她数算子女带来的安慰吧，但也要同样数算烦恼。她结婚时哭泣，含着眼泪许下多少誓愿。她怀孕，在生子之前，她的生育能力已带来烦恼。她生产时患病。多么甜蜜的抵押品，却始于危险又终于危险，带来痛苦先于快乐！它是以危险换来的，且不能随己意拥有。\par

26. 何必谈论哺乳、教养和婚嫁的烦恼？这些是那些幸运者的不幸。母亲有后嗣，但这只增加她的忧伤。因为我们不应谈论逆境，免得最圣善的父母胆战心惊。我的姐妹，请想想，那些不能说出口的事，承受起来该有多难。而且这还只是在今世。但日子将要来到，那时人们会说：\BibleQuote{不生育的，和从未生育过的胎，是有福的！}\BibleRef{路 23:29}因为今世的女子是被孕育，也孕育；但天国的女子却拒绝婚姻之乐和肉欲之乐，为能在身体和灵魂上都是圣的。\par

27. 我何必再谈妻子对丈夫所欠的辛苦伺候和服事呢？天主在奴隶出现之前，就已命令女人要服务。\BibleRef{创 3:16}我提及这些，是为了让她们更甘心情愿地顺服；她们的赏报，若蒙悦纳，便是爱；若不蒙悦纳，便是罪的惩罚。\par

由此生出种种恶欲的诱因，因为她们用各种颜色涂画面孔，不怕取悦自己的丈夫；从涂画面孔，进而想到玷污自己的贞洁。这是何等的疯狂，竟要改变天性的模样而寻求涂饰，又因畏惧丈夫的评判而放弃自己的本来面目。因为凡想要改变自己天然容貌的，就是第一个反对自己的人。她们努力取悦别人，却使自己不悦。女人啊，若论你的丑陋，还有什么比你害怕以真面目示人更真实的证据呢？你若美丽，为何隐藏自己？你若丑陋，为何假装美丽，以致既不满足于自己的良心，又得不到别人的错爱？因为他爱的是别人，你渴望取悦的也是别人。既然你本人教导他如此，他若爱上别人，你又何必生气呢？你成了伤害自己的邪恶教导者。\par

其次，即便是美丽的妻子，为求不令人厌倦，又需要多少花费呢？昂贵的项链挂在她的颈上，另一边，织金的袍子在地上拖曳。这究竟是买来的炫耀，还是真实的财富？又使用了多少种香料的诱惑！耳朵被宝石压得沉重，眼中滴入与天然不同的颜色。在这一切改变之下，还有什么属于她自己呢？已婚妇人珍爱自己的感受，难道她认为这就是生活？\par

可是你们，幸福的童贞女啊，你们不知道这些磨难——不如说是装饰——你们圣洁的端庄，流露于羞怯的面颊，甜美的贞洁便是你们的美貌；你们并不专注于男人的眼光，不以他人错误所得来的好处为功德。你们也有自己的美丽，由德行的优雅而非肉体的容貌所供给，岁月不能终结，死亡不能夺去，疾病也无法损害。只应寻求天主作为优美品德的评判者，祂甚至在不甚美丽的身体中，喜爱更美丽的灵魂。你们不知生育的重担与痛苦，但虔敬的灵魂所生的子女更多，它视众人为子女，后嗣兴旺，毫无丧亡，不知死亡，却有许多继承人。\par

所以，圣教会不知婚姻，却生养繁盛，在贞洁上是童贞女，在子女上却是母亲。她，这位童贞女，不是由人父，而是由圣神生了我们作她的子女。她生产我们，没有痛苦，却有天使的欢欣。她，这位童贞女，不是用肉身的奶，而是用宗徒的奶喂养我们，宗徒就是用这奶喂养那些尚为孩童的稚嫩子民的\BibleRef{格前 3:2}。因为哪一位新娘比圣教会子女更多呢？她在圣事中是童贞女，对子民却是母亲，甚至圣经也证实她的多产，说：\BibleQuote{因为荒胎者的子女，比有夫之妇的子女更多}？她没有丈夫，却有一位净配，因此，她——无论是万民中的教会，还是个体中的灵魂——毫不损伤端庄，将天主圣言作为永恒的净配，不受任何伤害，充满理性。\par

% structure: p0056 source=h2 target=subsection reason=relative-heading-rank; base=h2

\subsection{第七章}

\Emph{圣盎博罗削劝勉父母培养子女守童贞，并向他们陈明渴望孙辈所带来的烦恼。然而他也声明，并不禁止婚姻，反倒为婚姻辩护，反对那些反对婚姻的异端者。他虽将童贞置于婚姻之前，却也论及她们净配的美丽，以及净配用以装饰她们的恩赐，并将\BibleBook{}{雅歌}的某些经文应用于此。}\par

父母们，你们已经听到了，应以何种美德和追求来培养你们的女儿，使你们能拥有她们，借她们的功德赎去你们的过犯。童贞女是她母亲的奉献，借着她的每日牺牲，平息\OriginalTerm{天主的义怒}{divine power}。童贞女是父母不可分离的保证，她既不因嫁妆烦扰父母，也不离弃他们，更不以言行伤害他们。\par

然而，或许有人渴望得孙辈，并被称为祖父。首先，这样的人在追求他人的同时，放弃了自己的东西；在希望得到不确定的事物时，已失去了确定的；他舍弃自己的财富，却仍要被索取更多；若不付嫁妆，便被催逼；若他长寿，便成为负担。这是购买女婿，而不是赢得一个会把女儿卖给父母看一眼的女婿。难道她长久在母胎中孕育，就是为了转归他人权下吗？因此，父母忙于妆扮自己的童贞女，为的是让她更快地离开他们。\par

有人或许会说：那么，你是反对婚姻吗？不，我鼓励婚姻，并谴责那些惯于反对婚姻的人，甚至我惯于将撒辣、黎贝加、辣黑耳及旧约其他妇女的婚姻，作为卓越德行的实例。因为谁反对婚姻，就是反对生育子女，也就反对人类通过世代相继而延续的团体。因为若没有婚姻的恩赐激发生育的愿望，世代又怎能按继续的秩序代代相传呢？或者，又如何能说明依撒格因父亲的虔敬，作为祭品走上天主的祭台，或是在肉身中的以色列看见了天主\BibleRef{创 32:28}，并在反对他们赖以生存的婚姻时，给自己的人民起了一个圣名呢？那些人虽邪恶，但至少有一点，连明智者也认可，就是他们反对婚姻，正是承认他们本不该出生。\par

因此我并不反对婚姻，而是重述圣洁童贞的卓越。童贞只是少数人的恩赐，婚姻却是众人的。而且童贞本身若无某种产生的方式，也不可能存在。我将美善与美善相比较，为要显出哪一样更为卓越。我也不提出自己的任何意见，而是重复圣神借先知所说的话：\BibleQuote{不生育而又无玷的，是有福的。}\BibleRef{智 3:13}\par

36. 首先，在那些想要结婚的人最渴望的事上，就是夸耀她们丈夫的美丽，她们必须承认自己比不上童贞女，因为只有对她们才适合说：\BibleQuote{你比世人更美，你的嘴唇滴流恩宠。}那位净配是谁呢？是一位不沉迷于寻常享乐、不因拥有财富而骄傲的，而是祂的宝座永无穷尽。君王的女儿分享祂的荣耀：\BibleQuote{王后佩戴金饰，穿上各种美德织成的锦衣，侍立在你的右边。女儿！请听，请看，侧耳倾听：忘却你的民族和你父家！因为君王恋慕你的美艳，因为祂是你的天主。}\par

37. 你想一想，圣神藉着圣经的明证，赐给你的是何等国度——黄金与美丽。黄金，或是因为你是永恒君王的新娘，或是因为你心志不挠，不为享乐的诱惑所俘获，反而如王后一般统御它们。黄金，又因为这金属经过火炼更加宝贵，同样，祝圣于圣神的童贞身体，其荣美也愈发增长。因为谁能想象，那被君王宠爱、被审判者认可、献于主、祝圣于天主的女子，她的美丽有谁可及？她永为新娘，永不出嫁，如此，爱无止境，贞洁不损。\par

38. 这确是真美，无缺无欠，唯有这美配得听主说：\BibleQuote{我的爱卿，你全然美丽，毫无瑕疵。我的新娘，请随我离开\Place{黎巴嫩}，请随我离开\Place{黎巴嫩}！你从信仰之初，从\Place{色尼尔}和\Place{赫尔孟}山顶，从狮子的穴洞，从豹子的山岳，越过并穿过。}\BibleRef{歌 4:7} 藉此所表达的是，一个献于天主的祭台的童贞灵魂的完美无瑕之美。她不因可朽之物而动摇，身处灵性野兽出没的穴窟，却一心藉着天主的奥迹，使你堪当被爱者，祂的胸怀充满喜乐。因为\BibleQuote{美酒令人心神欢畅。}\par

39. \BibleQuote{你的衣香，}祂说，\BibleQuote{超越一切香料。}\BibleRef{歌 4:10} 又说：\BibleQuote{你衣服的芬芳，如同\Place{黎巴嫩}的馨香。}\BibleRef{歌 4:11} 童贞女啊，你看你所展示的进步。你第一种香气超越一切香料，那些香料是为安葬救主而用的，\BibleRef{若 19:39} 这香气源自克制肉身的冲动和肢体享乐的消亡。你第二种香气如同\Place{黎巴嫩}的馨香，发散出主肉身的不可朽坏，那是童贞纯洁的花朵。\par

% structure: p0066 source=h2 target=subsection reason=relative-heading-rank; base=h2

\subsection{第八章。}

\Emph{论及\BibleBook{}{雅歌}中关于蜂巢的段落，他加以阐述，将贞洁的童贞女比作蜜蜂。}\par

40. 那么，让你的工作如同蜂巢一般，因为童贞正可比作蜜蜂，它是如此勤劳、如此端庄、如此自制。蜜蜂以露水为食，不涉婚床，酿造蜂蜜。童贞女的甘露是天主的话，因为天主的言语如甘露降下。童贞女的端庄是无玷的本性。童贞女的出产是唇舌的果实，没有苦涩，充满甘甜。它们一起工作，果实也共同分享。\par

41. 我的女儿，我多么希望你效法这些蜜蜂：它们以花为食，它们的后代由口采集而聚合。我的女儿，效法它吧。不要让你的言语蒙上诡诈的面纱；不要有任何欺诈的掩饰，使它们纯洁并充满庄重。\par

42. 要让你的口结出永无穷尽的功劳。不要只为自己积蓄（因为你怎知道你的灵魂何时被索取？），免得你离开时，留下堆满谷物的粮仓，这粮仓对你今生和功劳都无帮助，而你却被催迫到那里，又无法携带你的财宝。所以，要向穷人富足，好使他们既然与你同享人性，也能同享你的财物。\par

43. 我也为你指出当采撷哪朵花，就是那一位说的：\BibleQuote{我是原野的花朵，谷中的百合，如荆棘中的百合。}这明显表明，美德被灵性邪恶的荆棘所包围，因此，若不谨慎靠近，无人能采得果实。\par

% structure: p0072 source=h2 target=subsection reason=relative-heading-rank; base=h2

\subsection{第九章。}

\Emph{讨论\BibleBook{}{雅歌}中其他与本题相关的段落，\Person{圣盎博罗削}劝勉童贞女寻求基督，指出何处可寻见祂。接着描述祂的完美，并将童贞女与天使相比。}\par

44. 所以，童贞女啊，你要振起圣神的双翼，飞越一切恶习之上，如果你愿达到基督：\BibleQuote{祂居高天，却垂顾卑微。}祂的仪态如\Place{黎巴嫩}的香柏，枝叶在云端，根柢在大地。因为它从天开始，在地终结，结的果实离天堂极近。你要殷勤寻找那珍贵之花，或许能在你心灵的深处找到，因为它常在卑微之处为人所享。\par

贞洁喜爱在园中生长，\Person{苏撒纳}在园中漫步时发现了它，并且宁愿死去，也不愿它受到玷污。但\Quote{园}的意思是由祂自己指出的，祂说：\BibleQuote{我的妹妹，我的新娘，是关闭的花园，是封锁的泉源；}因为在这种园中，纯净泉源的水映射出天主形象的容貌，以免其水流因与灵性野兽打滚之处的泥泞混合而受到污染。为此，贞女们的端庄，由圣神的壁垒圈起，被封闭起来，以免敞开而遭劫掠。于是，如同一个外人不可进入的园子，散发着紫罗兰的芬芳，洋溢着橄榄树的香气，焕发着玫瑰的光彩，好使虔诚在葡萄树中增长，和平在橄榄树中增长，而祝圣童贞的端庄在玫瑰中增长。这就是圣祖\Person{雅各伯}所闻到的香气，当时他听到父亲说：\BibleQuote{看，我儿的香气，犹如上主祝福的田地的香气。}\BibleRef{创 27:27} 因为尽管圣祖的田地几乎充满各种果实，但别的田地需以更大的劳苦产出庄稼，而这园子却产出花朵。\par

哦，贞女，如果你希望你的花园如此馨香，就当用先知们的规诫圈住它：\BibleQuote{要在你的口前设立守卫，在你的唇上安上门户，}好使你也能说：\BibleQuote{如同森林中的苹果树，我的爱人在男子中也是如此。在他的荫下，我欣然坐下，他的果实满口香甜。我找到了我灵魂所爱的，我拉住他不放。我的爱人下到自己的园中，去吃他树上的果实。我的爱人，我们到田间去吧。请将我有如印玺放在你心上，有如印玺放在你臂上。我的爱人皎白红润。} 因为，贞女啊，你应当完全知晓你所爱的那一位，并应在他身上认出他天主性体的一切奥秘和他所取的肉身。祂恰当地是洁白的，因为祂是父的光辉；又是红润的，因为祂由童贞女诞生。两种性体的颜色都在祂身上闪耀发光。但要记得，祂天主性体的标记比肉身的奥秘在祂身上更为古老，因为祂并非来自童贞女，而是已存在的那一位进入了童贞女。\par

祂，那位被兵士抢劫，被枪矛刺透，为要以他圣伤的血来医治我们的，必定会回答你（因为他是心谦良善、相貌温和的）：\BibleQuote{北风啊，醒来吧！南风啊，吹来！吹拂我的花园，使它的香气四溢。}因为从世界各地，圣善宗教的芬芳日益增长，祝圣童贞的四肢闪耀生辉。\BibleQuote{我的爱人，你美丽如\Place{提尔匝}，可爱如\Place{耶路撒冷}。}这样，那并非可朽肉身的美丽——它将随着疾病或衰老而终结——而是好德行的声誉，它不受任何灾难影响，永不消逝，这才是贞女们的美丽。\par

既然你已配得不再与人相比，而是与天上的存有相比——你在地上过着他们的生活——就请从主那里领受你要遵守的规诫：\BibleQuote{请将我有如印玺放在你心上，有如印玺放在你臂上；}为能彰显更明显的明达与行动的证据，在其中，基督——天主的肖像——能闪耀光芒，祂完全等同父的性体，表达着祂从父的天主性中所承受的一切。因此，宗徒\Person{保禄}也说，我们在圣神内受了印记；\BibleRef{弗 1:13}因为我们拥有子内的父的肖像，以及圣神内的子的印记。因此，我们既被这\OriginalTerm{天主圣三}{Trinity}所印记，就更应当勤奋谨慎，免得轻浮的性情或任何不忠的欺骗破坏了我们心中所领受的\OriginalTerm{质印}{pledge}。\par

但是，让敬畏为圣洁的贞女们确保这一切吧，教会首先为她们提供了这样的保护，她为柔嫩子女的繁荣而焦虑，自己就如城垣，双乳如同许多塔楼，增添自己的忧心，直到敌对攻击的恐惧结束，她以慈母的关爱为她强壮的子女赢得了和平。为此，先知说：\Quote{愿和平临于你的德能，愿丰足在你的塔楼里。}\par

随后，和平之主自己，以祂强壮的臂膀拥抱所托付给祂的葡萄园，眼见枝条发出嫩芽，便以欣喜的容色，为这些幼果调节微风，正如祂自己所证实的，说：\BibleQuote{我的葡萄园就在我眼前，\Person{撒罗满}有一千，看守果子的人有二百。}\BibleRef{歌 8:12}\par

上面说的是：\BibleQuote{六十位勇士环绕着它的子孙，手持利剑，精于战阵，}\BibleRef{歌 3:7}这里却有了一千二百。数目随着果实的增加而增加，因为一个人越圣善，就越受护守。因此，先知\Person{厄里叟}显示了那些临在护守他的天使军旅；因此，\Person{农}的儿子\Person{若苏厄}认出了天军统帅。那么，那些也能为我们作战的，就能护守我们内的果实。至于你们，圣洁的贞女们，有特殊的护守，因为你们以无玷的贞洁守护着主的净配的卧榻。如果那些以天使的生活方式作战的你们，有天使为你们作战，那也不足为奇。童贞的贞洁配得那些生命的护守，因为它已达到天使的生命。\par

我们为何还要用更多的话继续赞美贞洁呢？因为贞洁甚至造成了天使。谁持守了它，就是天使；谁失去了它，就是魔鬼。因此，宗教也由此得名：那作天主新娘的，是贞女；那为自己制造神祇的，是娼妓。关于复活，我又该说什么？你们已经拥有了它的赏报：\BibleQuote{因为在复活的时候，他们既不娶，也不嫁，}祂说，\BibleQuote{如同天上的天使。}\BibleRef{玛 22:30} 那应许给我们的，已经临在于你们；你们祈祷的对象，已在你们内；你们属于这个世界，却又不在这世界内。这个世代曾抓住你们，却未能留住你们。\par

53. 然而，这是何等伟大的事：天使因\OriginalTerm{无节制}{incontinentia}而从天堕落于尘世，贞女却因\OriginalTerm{贞洁}{castitas}由尘世升入天堂。有福的贞女，肉身的欢愉不能引诱她们，逸乐的污秽不能使她堕落。节制的饮食与戒酒使她们不染恶习，因为这样保守她们不知罪恶之根由。那引致罪恶的，常欺骗义人。如此，天主的子民坐下吃喝后，便否认了天主。\BibleRef{出 32:5} 如此，\Person{罗特}不知情，竟容忍了女儿的恶行。\BibleRef{创 19:32} 同样，\Person{诺厄}的儿子倒退着遮盖父亲的裸体，那放荡者看见了，端庄者却羞愧，并恭顺地遮住，唯恐自己若看见亦会冒犯。\BibleRef{创 9:22} 酒的力量何等大，以致洪水未能使他赤裸，酒却使他赤裸。\par

% structure: p0084 source=h2 target=subsection reason=relative-heading-rank; base=h2

\subsection{第十章}

\Emph{最后，提及\OriginalTerm{童贞}{virginitas}的另一光荣，即它免于贪婪。\Person{圣安博罗削}向他的姐妹讲话，提醒她那些免于奢侈和虚荣之烦恼的人是何等幸福，而这些烦恼会临到即将结婚的人身上。}\par

那么怎样呢？多么幸福，没有对财物的贪欲煽动你们！穷人索求你们所有的，却不要求你们所无的。你们劳苦的果实是穷乏者的宝藏，而两个小钱，若那是某人的所有，对施予者而言便是财富。\par

54. 那么，我的姐妹，请听，你避开了什么。因为我无须教导你，你也无须学习应提防什么，因为全德的实践不需教导，反教导他人。你看那犹如游行中的轿子，为取悦而精心打扮的人，她吸引所有人的注目凝视；她因竭力取悦反而不美，在她取悦丈夫之前，已使众人不悦。但在你身上，摒弃一切对华丽外表的用心更为相宜，你不装饰自己，这本身就是一种装饰。\par

55. 请看那被伤痕穿孔的耳朵，并怜悯那被重负压弯的颈项。金属的种类不同，并未减轻痛苦。一种情况是链子捆绑颈项，另一种是脚镣束缚双足。身体被黄金或铁所负，并无分别。如此，颈项被压弯，步履受阻碍。其价格并未使情况更好，除了你们女人害怕那使你们受苦的东西丢失外。他人的判决或你们自己的判决判你们有罪，有何区别？不，你们比那些人更悲惨，你们被公共司法定罪，因为他们渴望获释，你们却渴望被捆绑。\par

56. 然而多么可悲的处境：适婚的她，犹如待售之物，被置于拍卖中供人出价，以致出价最高者购得她。奴隶的出售条件更可容忍，因为他们常能选择主人；若少女选择自己的主人，便是冒犯，若不选择，则是侮辱。而她，即使美貌可爱，既恐惧又期望被人看见；她期望如此，以便将自己卖得更好价钱；她又害怕被人看见，这事实本身就不合宜。但那些愿望、恐惧和猜疑是何等荒谬，皆在于揣度求婚者将是何等样人：唯恐贫穷者欺哄她，或富有者蔑视她，唯恐俊美求婚者嘲笑她，唯恐高贵者轻看她。\par

% structure: p0090 source=h2 target=subsection reason=relative-heading-rank; base=h2

\subsection{第十一章}

\Emph{\Person{圣安博罗削}答复那些针对他赞颂\OriginalTerm{童贞}{virginitas}的劝告无用的异议，他提出贞女的实例，特别在他提及的各个地方，并讲述她们在该事业中的热忱。}\par

57. 有人或说，你总是在歌颂贞女。我既总是歌颂她们却无成效，又当如何？但这不是我的过错。而且，贞女从\Place{普拉森舍}，或从\Place{波诺尼亚}，以及\Place{毛里塔尼亚}前来受祝圣，为在此领受帕巾。你看，这里有一件引人注目的事。我在此处理此事，却说服了别处的人。若真如此，且让我在别处谈论这主题，好使我能说服你们。\par

58. 那么，这是怎么回事：连那些未听我讲的人遵循了我的教导，而那些听我讲的人却不跟从我？因为我认识许多贞女，她们有这愿望，却被她们的母亲阻止前行，而更严重的是，这些母亲是寡妇，我现在要对她们说话。因为如果你们的女儿渴望爱人，她们能按法律选择自己愿意的人。那么，那些被允许选择男人的女子，难道不被允许选择天主吗？\par

59. 请看，\OriginalTerm{端庄}{pudicitia}之果何等甘饴，它甚至生发在野蛮人的情感中。来自\Place{毛里塔尼亚}极远之地，这一边和那一边的贞女，渴望在此受祝圣；尽管所有家族都身处束缚之中，然而端庄不能被束缚。那哀叹为奴之苦的女子，宣认了一个永恒的国度。\par

60. 我该如何述说\Place{波诺尼亚}的贞女呢？她们是一群丰产的\OriginalTerm{贞洁}{castitas}队伍，她们弃绝世俗的欢愉，居住在\OriginalTerm{童贞}{virginitas}的圣所中。她们并非生活在世俗群体中的性别，在共同的贞洁内达到二十人的数目，并获得百倍的果实，她们离开父母的居所，涌进基督的家园，如同不懈贞洁的战士；她们时而歌唱灵性的歌曲，以劳作维生，并用自己的双手寻求资源以行慷慨。\par

61. 但若寻找贞女的吸引力增强了（因为她们比其他人更致力于搜寻和守望纯洁），她们便以极大的恒心，追踪隐藏的猎物，直至其内室；或者，若有谁的飞翔似乎更为自由，你会看见她们振翅而起，听见她们羽翼的窸窣声和欢呼的爆发；以至于用贞洁的端庄队伍，将那飞翔中的一人团团围住，直到她在美好的陪伴中欢欣，忘记父家，进入\OriginalTerm{端庄}{pudicitia}之境，及\OriginalTerm{贞洁}{castitas}的围护之家。\par

% structure: p0097 source=h2 target=subsection reason=relative-heading-rank; base=h2

\subsection{第十二章}

\Emph{父母鼓励对贞洁生活的向往是非常值得渴望的，但当对天主的爱吸引一位少女，甚至违背他们意愿时，更值得称赞。不应惧怕父母的暴力或财产的损失，\Person{圣盎博罗削}就讲述了一个这样的实例。}\par

那么，父母的热心，有如顺风，帮助贞女，这是好事；但更光荣的是，即使没有年长者的激励，年轻心灵自身的火焰也迸发出贞洁之火。父母会拒绝嫁妆，但你有一位富裕的净配，因祂的财富而满足，你不会失去父家产业的收益。为贞洁而贫穷，比嫁妆好多少！\par

然而，你听说过谁因为渴望贞洁而被剥夺了合法的遗产吗？父母反对她，却乐意被战胜。他们起初抗拒，因为不敢相信；他们常常生气，为让人学会战胜；他们威胁剥夺继承权，试探人是否不怕暂时的损失；他们用精致的诱惑来爱抚，看人是否不被各种享乐的诱因软化。啊，贞女，你被催促时，是在受锻炼。父母焦急的恳求是你的第一场战斗。啊，少女，先征服你的亲情。你若征服了家庭，便征服了世界。\par

但假设你祖业的丧失等待着你；未来的天界岂非对可朽脆弱财物的补偿？因为我们若相信那来自天上的信息，\BibleQuote{没有人为了天主的国，而舍弃房屋、父母、兄弟、妻子或儿女，而不在现世得到多倍的赏报，并在来世获得永生。}\BibleRef{路 18:29} 将你的信德托付给天主，祂曾将金钱托付给人；借贷给基督吧。你望德的寄托之忠实守护者，会以多倍的利息偿还你信德的\OriginalTerm{塔冷通}{talent}。真理不欺骗，正义不周旋，德能不虚诳。但如果你不信天主的话，至少该相信实例。\par

在我的记忆中，有一个在世俗中曾显贵、现在在天主眼中更高贵的少女，因父母和亲属催逼她结婚，便避居在圣洁的祭台前。一个贞女还能逃往哪里，比那奉献童贞祭献的地方更好呢？但她的勇敢远不止于此。她，那谦逊的奉献，贞洁的牺牲，正站在天主的祭台前，时而将司铎的右手放在头上，祈求他的祈祷，时而因义人延迟而不耐烦，便把头顶放在祭台下。她说道：\Quote{还有什么头纱，能比那祝圣诸头纱的祭台更好地遮盖我呢？这样的婚礼头纱最合适，因为万有的元首基督每天都在其上受祝圣。我的亲属们，你们在做什么？为什么还要以寻求婚姻来烦扰我的心？我早已为此作好了安排。你们给我提供一位新郎吗？我已找到更好的。你们尽可以展现我的财富，夸耀他的高贵，颂扬他的权能，我已拥有无人能相比的那位，于世间富足，于帝国掌权，于天庭尊贵。如果你们有这样一位，我并不拒绝这选择；如果找不到这样的，你们对我的好意，亲戚们，反倒成了伤害。}\par

当其他人沉默时，有一人颇为粗鲁地大声说：\Quote{如果你的父亲还在世，他岂能容许你不结婚？}然后她以更虔诚和更克制的敬爱之情回答说：\Quote{或许他正是离去了，好叫无人能阻碍我。}这回答论及她父亲，却成了对他自己的警告，因为他很快便以自己的速死证实了这话。于是其他人，各自害怕自己遭遇同样下场，便开始帮助她，不再像先前那样阻挠她。她的贞洁不但没有使她失去应得的产业，反而还获得了她完整操守的赏报。你们看，贞女们，热心所得的赏报；做父母的，要从这越轨的榜样中得到警戒。\par

\SourceRecord{Translated by H. de Romestin, E. de Romestin and H.T.F. Duckworth.；Nicene and Post-Nicene Fathers, Second Series；Vol. 10.；Edited by Philip Schaff and Henry Wace.；Buffalo, NY: Christian Literature Publishing Co.,；1896.；<http://www.newadvent.org/fathers/34071.htm>.}

% structure: p0001 source=h1 target=section reason=page-title

\section{论童贞（卷二）}

% structure: p0002 source=h2 target=subsection reason=relative-heading-rank; base=h2

\subsection{第一章}

\Emph{在本卷中，\Person{圣盎博罗削}打算以范例而非规诫来论述贞女的培育，并解释为何采用书面而非口头的方式。}\par

1. 在前一卷中，我曾渴望（尽管力不能及）阐明童贞之恩赐何其伟大，让这天上恩赐的恩宠本身吸引读者。在第二卷中，则宜教导贞女，并通过合宜的规诫来培育她。\par

2. 然而，鉴于我在劝导上软弱无力，又拙于教导（因为教导者理当胜过受教者），为避免看似放弃已承担的任务，或过于自负，我想，以范例而非规诫来教导更为可取；因为借助范例更能精进，盖因已经成就之事不视为困难，已经尝试之事被视为有益，而经由祖先德行的某种世代相传的实践传递给我们之事，在信仰上便具有约束力。\par

3. 但若有人指责我僭越，不如指责我热忱，因为我认为，对于那些向我提出请求的贞女，我甚至不应拒绝这一点。因为我宁可冒着损害自己谦逊的风险，也不愿未能成全她们的愿望，连我们的天主也以慈祥的嘉许眷顾她们的事业。\par

4. 我的工作也断不可被冠以僭越之名，因为她们既有可以学习的人，却仍寻求我的善意而非教导，且我的热忱或可原谅，因为当她们已有\OriginalTerm{殉道者}{martyr}的指引以恪守纪律时，我仍不认为，若我能将话语的劝说化为吸引她们发愿的诱惑，便是多余之举。善于教导者以严厉约束过失；而我，既不能教导，便以引诱为之。\par

5. 并且，由于许多不在场的人渴望能利用我的讲论，我便编纂了此书，好让他们手中持有所闻的实质内容，不至于以为，他们所持有的那人辜负了他们。但让我们继续我们的计划吧。\par

% structure: p0009 source=h2 target=subsection reason=relative-heading-rank; base=h2

\subsection{第二章}

\Emph{以\Person{玛利亚}的生平作为贞女的典范，详述她的诸多美德：贞洁、谦卑、艰苦的生活、喜爱独处等等；接着提及她对人的慈爱、学习的热情和常去圣殿的虔诚。而后，\Person{圣盎博罗削}描述她如何以这些美德为装饰，前来迎接无数贞女的队伍，并在凯旋中带领她们进入新郎的洞房。}\par

6. 因此，让\Person{玛利亚}的生平宛如\OriginalTerm{童贞}{virginity}本身，以一种肖像展现出来，犹如从明镜中，反映出贞洁的仪态与德行的形式。由此，你们可以获取生活的典范，它如范例般展示德行的明确准则：什么是你们必须纠正、成就和持守的。\par

7. 点燃学习热忱的首要因素，乃是师者的伟大。谁比\OriginalTerm{天主之母}{Mother of God}更伟大？谁比那被荣耀本身所选者更荣耀？谁比那未曾与另一身体接触而生育身体的她更贞洁？我又何必提及她的其他德行呢？她不仅在身体上，也在心灵上是\OriginalTerm{童贞女}{virgin}，从不以诡诈玷污心灵意向的纯真；她内心谦卑，言语庄重，心思明智，寡言少语，勤于阅读，不把希望寄托于无常的财富，而寄托于穷人的祈祷；专心工作，谈吐端庄；习惯寻求\OriginalTerm{天主}{God}而非人作为其思想的判官；不伤害任何人，友善对待一切人，在长者面前起立，不嫉妒同辈，避免自夸，遵循理性，热爱美德。她曾几何时，甚至以眼神使父母难过？她曾几何时，与邻居不和？她曾几何时，轻看卑微者？她曾几何时，回避贫乏者？她只去那种慈悲不以为耻、端庄不会回避的人群聚集。她眼中毫无阴沉，言语毫无轻率，举止毫无失礼，没有轻浮的动作，没有放肆的步伐，声音也不乖戾，以致她外在的仪表本身，便是她灵魂的肖像，是受赞许者的反映。因为一座井然有序的房屋，在门槛上就应当被认出，在初次进入时就应表明内中并无隐藏的黑暗，正如我们的灵魂不受身体的任何束缚阻碍，能如置放于内的灯盏般向外照耀。\par

8. 我何必详述她饮食的菲薄与服侍的繁多——一个超出天性地丰盛，另一个则几乎不足以维持天性？而且她从不懈怠，而是禁食的日子接连不断。即使有时渴望补充体力，她的食物通常只是随手可得之物，用以抵御死亡，而非谋求舒适。她是出于必需而非嗜好才安睡，然而身体沉睡时，她的灵魂却醒着，在睡眠中常常或重温读过的东西，或继续被睡眠打断的思绪，或实行所计划的，或预见将要实行的。\par

9. 她不惯于离开家门，除非是为敬礼，并且总有父母或亲戚陪伴。在家时私下忙碌，在外时有人陪同，但再没有比自己更好的监护者了，她以步态与举止令人起敬，与其说是以脚步前行，不如说是以德行一级一级地攀升。然而，虽然这位\OriginalTerm{童贞女}{Virgin}另有他人保护她的身体，她却独自守护自己的品格；她若作自己的导师，能学到许多，因为她具备一切美德的完美，因为她所做的一切都是教导。\Person{玛利亚}事事专注，仿佛有多人提醒她，她履行德行的每一项义务，仿佛她是在教导而非学习。\par

10. 圣史如此呈现了她，天使如此寻见了她，圣神如此拣选了她。何必详述细节呢？她的父母如何爱她，陌生人如何称赞她，她何等堪当天主之子由她诞生。天使进入时，她正独处家中，没有同伴，以免有人打断她的专注或搅扰她；她也不渴望任何妇女为伴，因她已有善念相伴。而且，当她独处时，她觉得自己并不那么孤单。她怎会孤单呢？她身边有那么多经书、那么多总领天使、那么多先知？\par

11. 同样，当\Person{加俾额尔}拜访她时，便这样见到她\BibleRef{路 1:28}；\Person{玛利亚}因看见人形而感到惊慌不安，但一听到他的名字，便认出他是她所认识的。因此，她对世人如同陌生人，但对天使却不然；由此我们可知道她的耳朵是端庄的，她的眼睛是含羞的。当被请安时，她保持沉默；当被问话时，她才回答。她起初心神不安，随后便许诺服从。\par

12. 圣经指出她对邻人是如何端庄自持的。因为当她得知自己被天主拣选时，变得更加谦卑，便立刻起身前往山区她的亲戚那里，这并不是为了借任何外在的事取得相信，因为她已相信了天主的话。她曾说：\Quote{你是有福的，因为你相信了。}她同那亲戚一起住了三个月。\BibleRef{路 1:56}在这段时间内，并非为寻求信德，而是表露善意。这事以后，那胎儿在母腹中欢跃，向主的母亲请安，在诞生前已具有了理性。\par

13. 随后，在许多奇事中：不生育的生了儿子，童贞女怀孕，哑巴说话，贤士来朝，\Person{西默盎}期待，星辰预示。\Person{玛利亚}曾因天使的进入而惊慌，却不为这些奇事所动。经上说：\BibleQuote{玛利亚把这一切事默存在自己心中。}\BibleRef{路 2:19}她虽是主的母亲，却仍切望学习主的诫命；她虽生了天主，却仍切望认识天主。\par

14. 她也每年在逾越节的庆日上\Place{耶路撒冷}去，与\Person{若瑟}同行。在童贞圣母身上，无论在何处，端庄总是她那独特美德的女伴。这是童贞不可或缺的伴侣，没有它，童贞便不能存在。因此，\Person{玛利亚}即使前往圣殿，也不缺端庄的护守。\par

15. 这便是童贞的肖像。因为\Person{玛利亚}的榜样本身便是万民的教训。因此，如果这作者不令我们反感，就让我们检验这作品，好使那渴望为自身赢得奖赏的人，效法这模范。多少美德在一位贞女身上闪耀！端庄的奥秘、信德的旗帜、虔诚的服侍——家中的贞女，服务的女伴，圣殿中的母亲。\par

16. 啊！她将迎接多少贞女，拥抱多少贞女，引领她们到主前，说：\BibleQuote{她忠于她的净配，忠于我的圣子；她以无玷的端庄守护了她的婚床。}主自己又将如何把她们举荐于父，重述祂的话说：\BibleQuote{圣父，这些是我为你保守的人，人子曾倚靠他们而安息；我祈求，我在哪里，他们也同我在一起。}\BibleRef{若 17:24}若这些不独为自己而生活的人，不仅应当造福自己，那么一位贞女可救赎她的父母，另一位可救赎她的兄弟。\BibleQuote{圣父，世界不认识我，这些人却认识我，他们决意不认识世界。}\BibleRef{若 17:25}\par

17. 那将是一个何等的行列，何等欢跃的天使们的喜悦，当那在地上度过天使般生活的人，堪当住在天堂时！那时，\Person{玛利亚}也拿起手鼓，激起贞女歌咏团，向上主歌唱，因为她们穿越了此世的海洋，却未受此世波浪的侵袭。\BibleRef{出 15:20}那时，每一位都将欢欣地说：\BibleQuote{我要走向天主的祭台，到使我青春喜乐的天主前；}以及\BibleQuote{我要向天主献上感恩，并向至高者偿还我的誓愿。}\par

18. 我也毫不犹豫地接纳你们到天主的祭台前，因为你们的灵魂，我毫不犹豫地称之为祭台，基督每日在其上为身体的救赎而被祭献。因为如果贞女的身体是天主的宫殿，那么她的灵魂，当身体的灰烬如被抖落，由永恒大司祭的手再度揭开时，会散发出神圣之火的馨香。有福的贞女们，你们借着天主的恩宠散发芬芳，如同花园通过花朵、圣殿通过宗教、祭台通过司铎一样。\par

% structure: p0024 source=h2 target=subsection reason=relative-heading-rank; base=h2

\subsection{第三章}

\Emph{\Person{圣盎博罗削}既已提出童贞\Person{玛利亚}为生活的榜样，便引证\Person{德格拉}为学习如何死去的模范。\Person{德格拉}并未受那些她曾被判投给的猛兽的伤害，反而从它们获得了敬畏的表示。他然后继续引入一个较近的实例。}\par

19. 那么，让圣母\Person{玛利亚}教导你生活的纪律，让\Person{德格拉}教导你如何被祭献，因为她躲避了婚姻的结合，并因丈夫的忿怒而被定罪，却因猛兽对童贞的敬畏，改变了它们的性情。因为她被预备给猛兽时，当躲避人的目光时，她将自己的要害部位献给一只凶猛的狮子，使那些原以无礼目光避开的人，转而端庄地注视。\par

20. 人们看见那野兽躺在地上，舔着她的脚，无声地表明它不能伤害贞女神圣的身体。这样，野兽敬重了他的猎物，忘记了自己的本性，穿上了人类所丧失的本性。人们似乎通过某种本性的转移，看到披着野蛮的人驱赶野兽施暴，而野兽却亲吻贞女的脚，教导人们何为应尽的本分。贞洁本身具有如此令人钦佩之处，连狮子也钦佩它。食物没有引诱它们，尽管它们未进食；冲动没有催促它们，尽管它们被激怒；愤怒没有使它们暴躁，尽管它们被煽动；习惯也没有盲目地驱使它们如往常一样；它们自己的天性也没有让它们充满凶残。它们在敬重殉道者时树立了虔敬的榜样；在只亲吻贞女的脚，眼睛转向地面，仿佛出于羞涩，生怕任何男性，甚至是野兽，看见贞女赤身露体时，给出了一堂赞颂贞洁的课。\par

21. 有人会说：\Quote{你为什么提出\Person{玛利亚}的榜样，好像有人能效法主的母亲呢？还有\Person{德格拉}的榜样，她是外邦人的宗徒所训练的那位？如果你想要门徒，就给我们和我们同类的一位教师吧。} 因此，我要把这类的一个近期榜样摆在你面前，好使你明白，那位宗徒不仅是某一个人的教师，也是众人的教师。\par

% structure: p0029 source=h2 target=subsection reason=relative-heading-rank; base=h2

\subsection{第四章}

\Emph{在\Place{安提约基雅}，有一位贞女因拒绝向偶像献祭，被判处送入淫窟，她与一位基督徒士兵交换衣服后，安然逃脱。随后士兵为此被定罪，她又返回，两人争夺殉道的奖赏，最终这奖赏都赐给了他们两人。}\par

22. 不久之前在\Place{安提约基雅}，有一位贞女尽量避免公开露面，但她越是躲避人们的目光，人们就越是渴望见到她。因为听闻而未目睹的美更令人向往，激情有两种动力：爱和了解——只要没有遇到不那么令人喜欢的事物；被认为令人喜欢的东西被认为更有价值，因为在这种情况下，眼睛不是通过察看作判断，而是被爱燃烧的心灵充满了渴望。于是，这位圣洁的贞女，为避免他们的情欲因渴望得到她而不断增长，便声明了她保全贞洁的意愿，这样就熄灭了那些邪恶之人的欲火，以致她不再被爱，反而被控告。\par

23. 于是迫害兴起。这位少女不知如何逃脱，又怕落入那些图谋她贞洁的人手中，便为英勇的德行预备好了自己的灵魂，她如此虔敬，不怕死亡；如此贞洁，期待死亡。她受冠冕的日子到了。众人的期望达到了顶点。少女被带上来，作了双重的宣认，即宗教和贞洁的宣认。但他们看到她宣认的坚定、对自己端庄的担忧、对酷刑的预备以及被人注视时的脸红，便开始考虑如何用贞洁来压倒她的宗教，这样，夺去了她最重要的东西之后，或许也能夺去她所剩下的东西。因此判决是：她要么献祭，要么被送入淫窟。他们就这样为他们复仇的神明，是以何种方式敬拜呢？以这种方式作出判决的人，自己又是如何生活的呢？\par

24. 这位贞女对自己的宗教毫不迟疑，却担忧自己的贞洁，开始思考：我该怎么办？今天，殉道的冠冕和贞洁的冠冕都吝于给我。但是，在否认贞洁之创始者的地方，贞女之名是不被承认的。一个呵护妓女的人怎能是贞女？一个爱慕奸夫的人怎能是贞女？一个寻找情人的人又怎能是贞女？拥有贞洁的心灵比拥有贞洁的身体更可取。两者皆好，若两者兼顾；若不能，则让我保持贞洁，不是对人，而是对天主。\Person{辣哈布}也曾是妓女，但她信了天主之后，便获得了救恩。\BibleRef{苏 2:9} \Person{友弟德}也打扮自己，为讨一个奸夫的欢心，但她做这事是为了宗教而不是情欲，无人视她为淫妇。\EditorialNote{友 10} 这个例子结局良好。因为如果那位将自己交付于宗教的人，既保全了贞洁，又保全了国家，那么或许我，通过保全我的宗教，也将保全我的贞洁。但如果友弟德把贞洁置于宗教之上，当国家沦丧时，她也会失去贞洁。\par

25. 如此，她受到这些榜样的教导，同时牢记主的话，主说：\BibleQuote{谁为我的缘故丧失自己的性命，必要获得性命，} \BibleRef{玛 10:39} 她哭泣，沉默不语，连奸夫也不要听见她的声音，她不选择对自己的端庄遭受损害，却拒绝了对基督的损害。请想一想，她连自己的言语都加以守护，岂能容忍自己的身体遭受不贞呢？\par

26. 我的言辞这一刻变得羞怯，害怕继续赞美或描述那所发生之事的一系列邪恶。贞女们，掩上你们的耳朵！天主的贞女被带进羞耻之屋。但现在，贞女们，张开你们的耳朵。基督的贞女可以被暴露于羞辱，却不能被玷污。在任何地方，她都是天主的贞女，天主的圣殿，淫窟不能损害贞洁，贞洁却除去那地方的污名。\par

一大群放荡的男人冲向那地方。圣洁的贞女们，请听这殉道者的奇迹，忘记那地方的名字。门在里面关着，鹰隼在外面叫嚷；有些人在争抢谁先攻击猎物。但她举手向天，仿佛来到祈祷之所，而非淫欲之地，说道：\Quote{基督啊，祢曾为贞洁的\Person{达尼尔}驯服凶猛的狮子，\BibleRef{达 6:22} 祢也能驯服那些凶暴的人心。火焰对希伯来儿童成了露珠，海水因祢的仁慈、而非自身的本性而竖立，为犹太人让路。\BibleRef{出 14:22} \Person{苏撒纳}俯身受刑，却胜过了那些奸淫的控告者；那亵渎祢殿宇恩赐的右手便枯萎了；如今祢的殿宇自身也被亵渎；求祢不容许这亵圣的乱伦，因祢不容许偷窃。愿祢的圣名此时再次得荣耀，因我这为羞辱而来的人，得以贞女之身离去！}\par

她刚祈祷完毕，看哪！一个貌似可怕战士的人闯了进来。那贞女在他面前战栗，而颤抖的百姓曾给他让路。但她没有忘记所读过的内容。\Quote{\Person{达尼尔}，}她说，\Quote{曾去看\Person{苏撒纳}受刑，并独自宣判她无罪，尽管众人已定她的罪。羊羔可能隐藏在这豺狼的外形之下。基督也有祂的士兵，祂是万军之主。\BibleRef{玛 26:53} 或者，或许是行刑者进来了。我的灵魂，不要怕，这类人造就殉道者。贞女啊！你的信德救了你。}\par

士兵对她说：\Quote{姊妹，我恳求你，不要怕。我一个兄弟，来到这里是为救生命，非为毁灭。救我吧，这样你自己也得救。我像奸夫一样进来，若你愿意，我将以殉道者之身出去。让我们交换衣裳，我的适合你，你的适合我，各自为基督。你的长袍将使我成为真正的士兵，我的将保全你的贞洁。你将穿戴得宜，我将更好地赤身，好让迫害者认出我来。拿这衣服隐藏女性身份，给我那件以祝圣我为殉道者。穿上这大氅来遮掩贞女的肢体，但保全她的端庄。戴上这头巾遮盖你的头发并掩藏容貌。进入娼院的人惯于脸红。你出去以后，切勿回头看，记起\Person{罗特}的妻子，\BibleRef{创 19:26} 她虽以贞洁的眼睛看了不洁之物，却因此失去本性。不要怕祭献有任何缺失。我要为你向天主献上祭品，你为我向基督献上士兵。你已服了贞洁的善役，其工价是永生；你有正义的胸甲，以灵性的武装保护身体，有信德的盾牌以抵御伤害，还有救恩的头盔，\BibleRef{弗 6:14} 因为在基督所在之处，才有我们救恩的保障，因为男人是女人的头，而基督是贞女的头。}\par

他一边说，一边脱下外氅。这衣服先前一直被疑为迫害者和奸夫之衣。贞女伸出颈项，士兵献上外氅。那是何等景象，何等恩宠的彰显，他们竟在娼院中争夺殉道！也当思量这身份，一个士兵与一个贞女，即本性气质相异的人，却因天主的仁慈而相同，为应验那句话：\BibleQuote{那时，狼将与羔羊共食。} \BibleRef{依 65:25} 看啊，羔羊与狼不仅共食，且一同被祭献。我还需多说什么呢？换上衣服后，那少女从陷阱中飞出，不再靠自己的双翅，因她已被灵性的双翼承载，（历代未曾见过的景象）她以贞女之身离开娼院，且是基督的贞女。\par

但那些用眼看却看不见的人，如同强盗抢掠，又如豺狼扑向羔羊，发怒咆哮。一个更为无耻的人走了进去。但当他用眼看清状况后，便说：这是怎么回事？进去时是个少女，现在却见到一个男人。这不是以鹿代少女的古老寓言，而是贞女真成了士兵。我曾听闻基督变水为酒，却未信；如今祂又开始改变性别了。趁我们还是原来的样子，我们赶紧离开吧。难道我也被改变了吗？我看到的情况与我以为的大不相同。我来到娼院，却见到一个保证。然而我离开时也被改变了，因为我进去时不洁，出来时却是贞洁的。\par

当此事传开后，因这样一位征服者理当获得冠冕，他因那\OriginalTerm{童贞女}{virgin}而被定罪，而那位童贞女也被捕；结果不仅一位童贞女，更是一位\OriginalTerm{殉道者}{martyr}从青楼中走出。据说那位少女奔向刑场，二人都求死。他说：\Quote{我被判死刑，这判决本该拘禁我，却释放了你。}她答道：\Quote{我选你作保人，不是要你受死刑，而是作我\OriginalTerm{贞洁}{chastity}的保证。若贞洁受侵犯，我的性别仍在；若要流血，我不要任何人替我担保，我自己有办法偿还。这判决原是对我下的，也是为我下的。无疑，若我让你作债务的担保，而我不在时法官将你的财产判给债主，你便会与我分担判决，我也会用我的财产偿付你的义务。若我拒绝，谁不会判我该受可耻之死？更何况我现在面对的是生死？让我清白地死，免得我有罪地死。在此事上无中庸之道；今日我若非背负你的血债，便是为自己的血作见证。若我速返，谁敢拒我于外？若我耽搁，谁敢宣我无罪？我欠律法更大的债，因我不仅自己逃跑，还导致他人之死。我的肢体既能承受死亡，就不能忍受侮辱。童贞女能接受创伤，却不能接受侮慢。我逃避的是耻辱，不是殉道。我把我的袍子给了你，但我并未改变我的\OriginalTerm{誓愿}{profession}。你若剥夺我的死亡，不是救我，而是诱使我。当心，求你，不要抗拒，不要冒险与我相争。不要夺走你给予我的恩惠。拒绝执行这判决，便是重新恢复前一个判决。因为判决已改回前一个。若后一个不约束我，前一个仍会约束。我们都可以满足判决，只要你让我先被处死。从你那里，他们不能索要其他惩罚，但她的贞洁在童贞女身上却面临危险。因此，你若让人看见你使一个淫妇成为殉道者，比你再使一个殉道者沦为淫妇更为光荣。}\par

你料想结局如何？二人相争，各获胜利，冠冕未曾分开，反成了两个。如此，这两位圣洁的殉道者，彼此施恩，一个引发了殉道的动机，另一个则实现了殉道的成果。\par

% structure: p0043 source=h2 target=subsection reason=relative-heading-rank; base=h2

\subsection{第五章}

\Emph{\Person{圣盎博罗削}叙述了毕达哥拉斯学派两位朋友\Person{达蒙}与\Person{皮西阿斯}的故事，并指出上一章所述之事更为可嘉。同时，将异教徒对待其神明而未受任何惩罚，与\Person{雅洛贝罕}的亵渎及其所受之惩罚加以对比。}\par

哲学家们的学派将\Person{达蒙}与\Person{皮西阿斯}——即毕达哥拉斯学派——推崇备至。其中一人被判死刑，请求给予时间以料理事务；狡诈的僭主不信能找到这样的人，便要求一个担保人，若那人延迟回来，担保人将代受刑罚。我不知道二人中哪一举动更为高尚。一人找到了担保人，另一人则自愿充当。于是，当被判刑者有所耽搁时，担保人面容镇静，不拒绝死亡。就在他被押赴刑场时，他的朋友回来了，引颈就斧。僭主惊讶于友情对哲学家而言比生命更可贵，便请求被他所判刑的两人接纳他为友。德行的恩惠如此伟大，竟感动了一位僭主。\par

这些事固然值得称赞，但与我们的例子相比则逊色。因为那两人是男子，而我们的例子中有一位童贞女，她必须先超越她的性别；那些人是朋友，这些人却互不相识；那些人将自己献给一位僭主，这些人却献给众多僭主；且这些僭主更为残忍，因为在前一事例中，僭主宽免了他们，在后一事例中，僭主却杀害了他们；在前者中，一人为形势所迫，在后者中，各人的意志都自由。再者，后者更为明智，因为对前者而言，他们热忱的终点是友情的乐趣，对后者则是殉道的荣冠，前者为世人奋斗，后者为天主奋斗。\par

既然我们提到了那个被判刑的人，那么也当谈谈他对他的神明的看法，好让你们判断，那些被自己的信徒所嘲笑的神明是何等软弱。因为他来到\Person{朱庇特}的神庙，命人将他神像头上所戴的金冠冕取下，换上羊毛冠冕，说金冠冕冬天冰冷，夏天沉重。他如此嘲笑他的神明，不能承受重量或寒冷。他同样看到\Person{阿斯克勒庇俄斯}的金胡须，便命人取下，说父亲没有胡须，儿子有胡须不合宜。他还从手持金碗的神像手中夺走金碗，说他理应领受神明所赐之物。他说，人们祈求神明赐下美物，而金是最美的；然而，若金为恶物，神明就不应拥有；若为善物，不如让懂得如何使用的人拥有它。\par

他们如此可笑，以致\Person{朱庇特}不能保护自己的衣袍，\Person{阿斯克勒庇俄斯}不能保护自己的胡须，因为\Person{阿波罗}尚未开始长胡须；所有那些被奉为神明者，也都不能守住他们手持的金碗，他们并非害怕被控盗窃，而是毫无知觉。那么，谁愿朝拜他们？他们既不能像神明那样护卫自己，也不能像人那样隐藏自己。\par

然而，在我们天主的圣殿里，邪恶的君王\Person{雅洛贝罕}取走了他父亲存放的礼品，并在圣祭台上献给偶像，他伸出的右手岂不立刻枯萎了？他所呼求的偶像岂能帮助他？于是，他转向主，祈求宽恕，他那因\OriginalTerm{亵圣}{sacrilege}而枯萎的手，立刻因真宗教而痊愈。如此，在同一个人身上显明了天主的仁慈与义怒的完整范例：当他献祭时，突然失去右手；而当他忏悔时，又获得了宽恕。\par

% structure: p0050 source=h2 target=subsection reason=relative-heading-rank; base=h2

\subsection{第六章}

\Emph{圣安博罗削在结束第二卷时，将书中可能有的任何善归功于贞女们的功德，并阐明在制定任何严厉规诫之前，应先以榜样鼓励她们，这在人间的教导和圣经中都是如此做的。}\par

39. 我，担任主教尚未满三年，虽无亲身经验教导，却从你们的生活方式中学到了许多，为你们，圣洁的贞女们，准备了这份献礼。因为在这么短的入教时间内，能有何经验可言？若你们在其中发现任何花朵，请将它们收集在你们生活的怀中。这些并非为贞女们制定的规诫，而是取自贞女们的实例。我的言辞描绘了你们德行的肖像，你们可以在这篇讲话的明镜中，仿佛看到自己庄重的映象。若你们因我的能力而感到任何喜乐，这本书的全部馨香都属于你们。既然人各有见，若我的论述中有何简朴之处，就让众人阅读；若有更刚毅之处，就让更成熟者加以验证；若有端庄之处，就让它贴近胸怀，染红面颊；若有如花之处，就让青春绽放的年纪勿轻视它。\par

40. 我们应当激起新娘的爱，因为经上记载：\BibleQuote{你当爱上主你的天主。}\BibleRef{申 6:5} 在婚宴上，我们至少当用祈祷的装饰来妆点头发，因为经上记载：\BibleQuote{拍掌顿足。}\BibleRef{则 21:14} 我们应当向那永不止息的婚宴撒播玫瑰。即便在这些尘世的婚姻中，新娘在接受命令之前，先受到欢呼迎接，免得在因仁慈而滋养的爱情坚固之前，严苛的命令伤害了她。\par

41. 马要学会喜爱轻拍颈项的声音，以免拒绝轭具，它们先以诱哄的言语接受训练，然后才受到管教的鞭打。但当马将颈项置于轭下后，缰绳收紧，马刺催迫，同伴牵引，御者号令。同样，我们的贞女也应首先以虔敬的爱嬉戏，在婚姻的门槛上赏叹天上婚床的金色支柱，看到门框饰以花叶的彩环，品尝内里乐师奏乐的欢愉；好使她不因恐惧而在服从主的召唤之前，从主的轭下退缩。\par

42. \BibleQuote{你从\Place{黎巴嫩}下来，我的爱卿！你从\Place{黎巴嫩}下来。}\BibleRef{歌 4:8} 我们应时常重复这节经文，好让她至少因受到主言语的呼唤而跟随，若有人不信人的言语。这能力并非我们自己造就，而是领受的；这是奥秘之歌的天上教导：\BibleQuote{愿他以口唇亲吻我！你的爱抚胜过美酒，你的香气芬芳怡人，你的名香四射。}\BibleRef{歌 1:2} 那整个欢乐之所发出嬉戏之声，激起赞许，唤起爱情。经文继续道：\BibleQuote{故此少女都爱慕你。愿你拉着我，让我们以你为我们的喜乐和欢欣；君王引我进入了内室。}\BibleRef{歌 1:3} 她从亲吻开始，从而进入了内室。\par

43. 如今她如此坚忍劳苦，精修德行，以至亲手开启门闩，走向田野，居于壁垒之中，而之初她却曾追随香膏的芬芳；当她进入内室后，香膏就转变了。且看她往何处去：经上说：\BibleQuote{若是墙，我们就在上面建造银城垛。}\BibleRef{歌 8:9} 那曾以亲吻嬉戏的她，如今建造塔楼，好使她不仅因诸圣德的高贵城墙环绕，以致敌人的攻击徒劳无功，更能建立圣善功绩的坚固防御。\par

\SourceRecord{Translated by H. de Romestin, E. de Romestin and H.T.F. Duckworth.；Nicene and Post-Nicene Fathers, Second Series；Vol. 10.；Edited by Philip Schaff and Henry Wace.；Buffalo, NY: Christian Literature Publishing Co.,；1896.；<http://www.newadvent.org/fathers/34072.htm>.}

% structure: p0001 source=h1 target=section reason=page-title

\section{论童贞（卷三）}

% structure: p0002 source=h2 target=subsection reason=relative-heading-rank; base=h2

\subsection{第一章}

\Emph{圣盎博罗削在此转而重提利伯里乌斯向马塞利纳授予贞女帕时的那篇训诲。他述及众人蜂拥而至那位饱餍所有人的新郎的婚宴，接着论及她在基督由童贞女诞生之日宣发圣愿是何等合宜，最后以热切的劝谕敦促她爱慕基督。}\par

1. 由于我在前两卷中论述有所游离，现在，圣善的姊妹，是时候重温那些备受怀念的\Person{利伯里乌斯}的训诲了，你曾常与我谈论这些训诲，因为一个人的圣德越高，他的话语也越加令人愉悦。因为当你在救主圣诞的日子，在\Place{圣伯多禄教堂}，通过更换服饰宣发了你的守贞圣愿时（还有什么日子能比那童贞女获得她婴孩的日子更好呢？），当时许多贞女环立，争相要成为你的同伴。\Quote{你，}他说道，\Quote{我的女儿，你渴求了一个美好的婚约。你看，为你净配的诞辰，聚集了多少群众，却没有一个人得不到食物而离去的。这就是那一位，祂在被请赴婚宴时，把水变成了酒。\BibleRef{若 2:9} 祂也要将纯洁的童贞奥秘赐给你，你曾隶属于物质的卑下元素。这就是那一位，祂在旷野里用五个饼和两条鱼饱了四千人。\BibleRef{路 9:13} 祂本可饱餍更多人；若当时有多人需要餍足，他们就会得到餍足。如今祂邀请了众人来赴你的婚宴，但现在供应的不再是大麦饼，而是从天而来的身体。}\par

2. 今天，祂确实按人的方式，由童贞女诞生，但在万有之前，由圣父所生，身体像祂的母亲，能力像祂的父亲。在地上是独生子，在天上也是独生子。出自天主的天主，由童贞女所生，来自圣父的正义，出自大能者的大能，出自光明的光明，与圣父平等；在能力上不分离，不因圣言的延伸或扩展而混乱，好像与父混合，而是因着生发而有别于圣父。祂是你的兄弟，\BibleRef{歌 5:1} 没有祂，天上的、海里的和地上的一切都不能存在。圣父的美善圣言，正如经上说：\BibleQuote{在起初}\BibleRef{若 1:1}，这里你看到了祂的永恒。经上又说：\BibleQuote{圣言与天主同在。}\BibleRef{若 1:1} 这里你看到了祂的能力，与父不可分、不可离。\BibleQuote{圣言就是天主。}\BibleRef{若 1:1} 这里你看到了祂非受生的天主性，因为你的信德必须从这相互关系中引出来。\par

3. 我的女儿，你要爱祂，因为祂是美善的。因为，\BibleQuote{除了天主一个外，没有谁是善的。}\BibleRef{路 18:19} 既然毫无疑问，子是天主，天主是美善的，那么也就毫无疑问，天主子是美善的。我说，你要爱祂。祂就是圣父在晨星之前所生的，因为是永恒的，祂如子从母胎中出来；祂如圣言从心中发出。祂就是圣父所喜悦的；\BibleRef{玛 17:5} 祂是圣父的手臂，因为祂是万物的创造者，又是圣父的智慧\BibleRef{格前 1:30}，因为祂由天主口中发出；\BibleRef{智 24:3} 是圣父的大能，因为天主性的圆满有形有体地居住在祂内。\BibleRef{哥 2:9} 圣父如此爱祂，将祂抱在怀中，放在自己的右边，好使你学习祂的智慧，认识祂的能力。\par

4. 那么，如果基督是天主的大能，天主曾几何时没有大能吗？圣父曾几何时没有圣子吗？如果圣父确实永远存在，那么圣子也必然永远存在。所以祂是完美父的完美子。因为凡削弱大能的，就是削弱那大能所属的祂。天主性的完美不容许不平等。所以，你要爱圣父所爱的，尊敬圣父所尊敬的，因为\BibleQuote{不尊敬子的，就是不尊敬父}\BibleRef{若 5:23}，且\BibleQuote{凡否认子的，就没有父。}\BibleRef{若一 2:23} 关于信德就说这么多。\par

% structure: p0008 source=h2 target=subsection reason=relative-heading-rank; base=h2

\subsection{第二章}

\Emph{接下来论及贞女的训练，他谈到饮食的节制，以及对心灵冲动的约束，引入了一些关于\Person{希波吕托斯}死亡与复活寓言的教导，并劝人戒避某些肉类。}\par

5. 但有时即使信德可靠，青春也不被信任。因此，要少用酒，以免身体的软弱加重，不可为了愉悦的兴奋，因为酒和青春，二者皆能燃起火焰。也要让禁食给稚嫩的年纪戴上辔头，让简单的饮食以某种缰绳约束尚未制伏的欲望。让理智加以阻止，望德加以克服，敬畏加以抑制。因为不认识如何管理自己欲望的人，就像被野马拖走的人一样，会被摔伤、擦伤、撕裂和致残。\par

6. 据说，这事发生在一个青年身上，因为他爱慕\Person{狄安娜}。但这寓言被诗人的故事渲染，说\Person{尼普顿}因对手被偏爱而心生忧伤，将疯狂送到了他的马身上，以此显示他的大能，因为他不是靠力气，而是靠诡计胜过了那青年。由于这件事，每年都为\Person{狄安娜}举行一次祭献，在祭台上献上一匹马。他们还说她是贞女，却（连娼妓都会为此感到羞耻）能爱上那不爱她的人。但就我而言，就让他们的寓言去发挥吧，因为尽管二者都是恶事，但一个青年如此迷恋一个淫妇以致丧命，比他们所说的两位天神争着犯奸淫，而\Person{朱庇特}又为他那犯邪淫的女儿报了仇，报复了那位治好那个曾在树林中侵犯\Person{狄安娜}的人之伤的医生，她无疑是最优秀的猎手，不是猎野兽，而是猎情欲：但也猎野兽，因此她赤裸着受崇拜，这还算较轻的恶了。\par

7. 那么，让他们把制伏疯狂的能力归于\Person{纳普敦}，好把不洁之爱的罪归于他。让他们把统治树林的权力归于\Person{狄阿娜}，她曾居住其中，好证实她所犯的奸淫。让他们把复活死人的能力归于\Person{埃斯库拉庇乌斯}，只要他们承认他遭雷击时自己未能逃脱。也让他们把并不属于他的雷电归于\Person{朱庇特}，好使他们见证他所承受的耻辱。\par

8. 我想，人应少食一切令肢体发热的食物，因为肉身拖累，甚至飞翔的鹰也被压下。但在你内，让那只我们所读到过的鸟，\BibleQuote{你的青春将更新如鹰}，在高处飞行，保持童贞的疾飞，不知对不必要食物的欲望。宴饮和应酬的聚集必须避免。\par

% structure: p0014 source=h2 target=subsection reason=relative-heading-rank; base=h2

\subsection{第三章}

\Emph{劝勉贞女们避见访客，持守端庄，效法\Person{玛利亚}的榜样，在奥迹的庆祝中保持静默。然后，在叙述了一个异教青年的故事和一位诗人的话之后，圣\Person{安博罗削}讲述了一位圣善司铎所行的奇迹。}\par

9. 我还愿意，年轻人之间的拜望，除了对父母及同龄人应有的以外，应当稀少。因为端庄因交往而磨损，大胆爆发，嬉笑潜入，害羞减弱，同时却讲究礼貌。不回答提问是幼稚，回答则是胡言。因此，我宁愿贞女少说话，而非多言。因为，如果妇女在教会中受命缄默，连神圣的事也是如此，且应在家问自己的丈夫，那么，我们应想，贞女当如何谨慎？她们的年纪因端庄而增色，静默更彰显她们的端庄。\par

10. \Person{黎贝加}来嫁给\Person{依撒格}，看见新郎，便取帕蒙上脸\BibleRef{创 24:65}，免得在结合之前让人看见，这岂是端德的微小标记？诚然，这美丽的贞女不是为了自己的美貌而害怕，而是为了保持端庄。至于\Person{辣黑耳}，当\Person{雅各伯}亲吻她后，\BibleRef{创 29:11}她哭泣呻吟，若非认出他是亲戚，不会停止哭泣。如此，她既守住了端庄所要求的，也不乏亲情。但若对男人说：\BibleQuote{不要注视少女，免得她使你跌倒}\BibleRef{德 9:5}，那么对一位已祝圣的贞女又该说什么呢？她若爱人，在心灵上犯罪；若被人爱，则在行为上也犯罪。\par

11. 静默之德，尤其在教会中，甚为伟大。不要让任何一句神圣训诲的话从你耳边溜走；如果你侧耳倾听，就要抑制自己的声音，不说一句想要收回的话，而应少发大胆之言。因为实在，在多言中难免有罪过。\BibleRef{箴 10:19} 对杀人者说：\BibleQuote{你犯了罪，静默吧！}\BibleRef{创 4:7} 免得他再犯罪；但对贞女必须说：\Quote{静默，免得你犯罪。} 因为，据我们读到，\Person{玛利亚}将一切关于她圣子的话都存记在心\BibleRef{路 2:19}，而你，当诵读到预告\Person{基督}将要来临，或显明祂已来临的章节时，不要以谈话制造喧噪，而要专注。还有什么比神圣的话语被谈话所淹没，以致听不见、不能相信或传扬，圣事因嘈杂的人声而模糊不清，为众人得救而献上的祈祷受到阻碍，更不合适呢？\par

12. 外邦人以静默向他们的偶像表示尊敬，兹举一例：当\Place{马其顿}王\Person{亚历山大}献祭时，一个为他点灯的外族童仆的衣袖着火，烧着了身体，他却纹丝不动，既不以呻吟流露痛苦，也不以静默的眼泪显出征象。一个外族童仆竟有如此敬畏的纪律，以致天性都被征服。然而他惧怕的不是那些并非神明的神，而是君王。因为他何须惧怕那些若被同样的火烧着也会燃烧的偶像呢？\par

13. 还有更好的例子：一个青年在父亲的宴席上受命不可用粗鄙的举止泄露其不洁之爱。那么，圣善的贞女，你在奥迹中要戒除呻吟、喊叫、咳嗽和嬉笑。难道你在奥迹中不能做到他在宴席上所做的吗？让童贞首先由声音显明，让端庄闭口，让宗教消除软弱，让习惯训导天性。让她的庄重首先向我宣示她是一位贞女：端庄的举止，稳重的步伐，羞怯的面容，让完整的凭证作为德行之军的先导。那见了以后还需询问的贞女，并不足以完全称许。\par

14. 有一个普遍的故事：当阵阵蛙鸣震响信众耳际时，天主的司铎命它们静默，并对神圣的话语表示敬意，于是喧声立刻平息。那么，沼泽可以静默而青蛙却不能吗？无理性的动物尚且能藉由敬畏而重新认出它们本性所不知的事吗？然而人却如此无耻，以至许多人对心灵宗教情感的尊重，竟不及他们对耳朵愉悦的在意。\par

% structure: p0022 source=h2 target=subsection reason=relative-heading-rank; base=h2

\subsection{第四章}

\Emph{总结了教宗\Person{利伯略}的劝导之后，圣\Person{安博罗削}继而讲述他姊妹的德行，尤其是她的禁食，不过他劝她在这方面要适度，并引证例子，要她在其他事上锻炼自己。他特别推荐天主经、夜间反复诵念圣咏，以及天明前背诵信经。}\par

15. 圣善记忆的\Person{利贝里乌斯}曾以类似方式劝勉你们，其言辞超越大多数人的实践，却不及你们的成就；你们不仅以德行获得了全备的操练，更以热忱超越之。我们受命实行禁食，但仅限于单日；你们却累积昼夜，度过无数无食的时光，若有人请求你们稍进饮食，放下书本片刻，你们立刻回答：\BibleQuote{人生活不只靠饼，而也靠天主口中所发的一切言语}。\BibleRef{玛 4:4} 你们的日常饭食仅由手边的食物构成，以致禁食比吃厌恶之物更可取；你们的饮品来自泉源，你们的哭泣与祈祷交织，你们的睡眠在书本上。\par

16. 这些君王适合更年少的人，而他以灰白的头发日渐成熟；但当一个贞女已征服自己的身体获得胜利时，应减少辛劳，以便保存为更年少者的教师。结实累累的葡萄藤若不时加抑制，很快就会折断。但当其尚幼时，任其茂盛生长，待其渐老则修剪，免得长成枝桠丛生，或因过度结实而丧命。一位好农夫通过照料土壤，使葡萄树保持良好状态，保护它免受寒冷，并防止正午的日头把它晒焦。他轮番耕作他的土地，如若不让它休耕，便交替种植庄稼，使田地通过作物轮换得以歇息。身为守贞的老练者，你也当如此，至少要为你心田播下不同的种子：时而适度进食，时而施行节俭的禁食，还有阅读、工作和祈祷，使劳苦的变换如同休憩的停战。\par

17. 整片土地并不出产同样的收成。这边山上生长葡萄，那边可见紫橄榄，别处有芬芳的玫瑰。强壮农夫放下犁头后，用手指刨土以种植花卉的根株，还用他转动葡萄藤间奋力耕牛的粗糙双手，轻柔地挤着羊的乳房。土地越是果实累累，便越是好地。因此，你当效法好农夫的榜样，避免用永久的禁食如同深耕般撕裂你的心田。让贞洁的玫瑰在你的园中绽放，还有心灵的百合，让紫罗兰圃从圣血的源头畅饮。有一句俗语说：\Quote{你想做得丰盛的事，有时全然不要做}。在四旬期的日子里总该有所增添，但不可出于炫耀，而应出于虔诚。\par

18. 频频祈祷也把我们举荐于天主。先知曾说：\BibleQuote{我每日七次赞美你}，尽管他忙于王国的事务，那么，诵读\BibleQuote{醒寤祈祷罢！免陷于诱惑}的我们，又该当如何呢？\BibleRef{玛 26:41} 当然，我们惯常的祈祷当伴随感恩，在起床时、出门时、准备进食时、进食后、上香时分，以及最后即将安歇时，都应诵念。\par

19. 我还要劝你在寝室内，无论是醒来时，还是在睡眠沾润身体之前，将圣咏与主祷文频频交替唱诵，好使你在休息伊始，睡眠发现你已脱去世俗的挂虑，正默想着天主的事。确实，最初发明\OriginalTerm{哲学}{Philosophy}这一名称的人，每天就寝前，令吹笛者演奏较柔和的旋律，以抚慰他为世俗烦扰的心。但他有如一个洗瓦片的人，徒然渴望以世俗的方法驱除世俗的事，因为他实在是给自己涂上新的泥浆，寻求享乐中的赏报；而我们却当擦净尘世诸恶的污秽，从肉身的每一点染污中洁净我们至深的灵魂。\par

20. 我们也当特别地在每日天亮前，诵念\OriginalTerm{信经}{Creed}，犹如盖在心上的印记，每逢惧怕何事，就在思想中重温它。因为营帐中的士兵或战场上的武士，何时没有他的军人誓言呢？\par

% structure: p0030 source=h2 target=subsection reason=relative-heading-rank; base=h2

\subsection{第5章}

\Emph{圣安博论及眼泪，解释达味的言语：\BibleQuote{每夜我将以我的眼泪浸湿我的床榻}，并继续讲论基督承担我们的忧苦和病弱。凡事都应归光荣于他，我们当以属灵的喜乐欢欣，而非世俗的方式。}\par

21. 现在谁还能不明白圣先知是为教训我们而说：\BibleQuote{每夜我将浸湿我的床榻，以我的眼泪浇灌我的床}呢？如果你按字面理解他的床，他表明应流出如此丰沛的眼泪，以致浸湿床榻，以泪浇灌祈祷者的卧榻；因为哭泣关乎现世，赏报在于将来，如经上所说：\BibleQuote{你们哭泣的是有福的，因为你们将要欢笑}；\BibleRef{路 6:21} 或者，如果我们把这先知的话应用于我们的身体，我们便应以忏悔的眼泪洗去肉身的过犯。因为\BibleQuote{撒罗满用黎巴嫩木为自己造了一架床，床柱是银的，床底是金的，床背镶嵌宝石。}\BibleRef{歌 3:6} 这床除了是我们身体的模型，还会是什么呢？因为宝石代表空气光辉的灿烂，火由金代表，水由银代表，地由木代表，人的身体就是由这四种元素构成，我们的灵魂安息其中，只要它不因山岭的崎岖或潮湿的地面而失去安息，而是被高举于罪恶之上，由这木头所支撑。为此达味也说：\BibleQuote{上主必在他病痛的床上赐他帮助}。怎么能说那不能感受痛苦、没有知觉的床是病床呢？但痛苦的身体就像那死亡的身体，正如经上所说：\BibleQuote{我这个人真不幸啊！谁能救我脱离这取死的身体呢？}\BibleRef{罗 7:24}\par

22. 既然我插入了一个提及主的身体的句子，免得有人因读到主取了受苦的身体而感到不安，他应记得主曾为\Person{拉匝禄}的死忧伤哭泣，\BibleRef{若 11:35}，并在祂的苦难中受了伤，从伤口流出了血和水，\BibleRef{若 19:34}，且交付了祂的灵魂。水为洗涤，血为饮用，圣神为祂的复活。因为基督独自为我们是望德、信德和爱德——望德在于祂的复活，信德在于圣洗，爱德在于圣事。\par

23. 正如祂取了受苦的身体，同样在软弱中祂也翻转了病榻，因为祂将其转化为对人肉体的益处。因为因祂的苦难，软弱终结了；因祂的复活，死亡也被终结。然而你们应当为世俗哀伤，却要在主内喜乐；为补赎而忧愁，却要为恩宠而喜乐。尽管外邦人的师傅也以有益健康的规诫，吩咐我们要与哭泣的一同哭泣，与喜乐的一同喜乐。\BibleRef{罗 12:15}\par

24. 但凡愿意解决这问题所有难题的人，应求助于同一位宗徒。他说：\BibleQuote{你们无论作什么，在言语或行为上，都该因主耶稣的名，借著祂感谢天主圣父。}\BibleRef{哥 3:17}因此，我们应将一切言行都归于基督，祂从死亡中带出生命，从黑暗中造出光明。正如一个生病的身体，一时用温暖抚慰，另一时用清凉舒缓，若按医师的指导变换疗法，便有益健康；但若违背他的吩咐，则加重病情；同样，凡是奉献给基督的，都是良药；凡是凭我们自己意愿所做的，都是有害的。\par

25. 因此，心灵的喜乐应当是意识到正义，并非因放纵的宴席或婚礼的音乐而激动，因为在这样的场合，端庄并不安全，而过度跳舞伴随欢庆时，诱惑便可能潜伏。我愿天主的贞女们远离这些。正如这世俗的一位师傅所说：\Quote{清醒的人不会跳舞，除非他疯了。}如今，若按此世的智慧，醉酒或疯狂是跳舞的原因，那么圣经中所记载的事例给了我们何等警示：\Person{若翰}，基督的前驱，因一个跳舞女子的愿望而被斩首，这例证表明跳舞的诱惑比亵圣的愤怒更为有害。\par

% structure: p0037 source=h2 target=subsection reason=relative-heading-rank; base=h2

\subsection{第六章}

\Emph{提到洗者后，\Person{圣盎博罗削}开始描述有关他死亡的事件，并反对跳舞和恶人的欢宴。}\par

26. 既然我们不应草率地提及如此伟大的人物，让我们细思他是谁，被谁，出于什么原因，怎样，以及在何时被杀。他是个义人，却被奸淫者处死，而一个死罪的刑罚，竟由有罪者转嫁于判官。再者，跳舞女子的赏报是先知的死亡。最后（这对所有蛮族都是一种荣誉），残酷的判决在欢宴与庆祝之中宣布，致命罪行的消息从宴席带到监牢，又从监牢带回宴席。一个邪恶的行为里包含了多少罪行！\par

27. 一场死亡之宴以王室的奢华摆设着，当比平时更多的人聚集时，皇后的女儿被从内室召来，带到众人面前跳舞。她从奸淫者那里能学到什么，除了丧失端庄？有什么比用不雅的舞姿，裸露那些本性隐藏或习俗遮掩的身体部位，以眼目嬉戏，扭动颈项，松散头发，更助长情欲呢？下一步恰是冒犯天主。因为在有跳舞、喧闹和击掌的地方，还能有什么端庄呢？\par

28. \BibleQuote{那时，王欢喜，对少女说：\InnerQuote{你无论向我求什么，我都必给你。}又发誓说：\InnerQuote{你无论求什么，即使是我王国的一半，我也必定给你。}}\BibleRef{谷 6:22}看，世俗的人如何判断他们的世俗权力，竟为了跳舞而给予王国。但那少女受她母亲的唆使，要求把\Person{若翰}的头放在盘子里带给她。所谓\BibleQuote{君王便忧郁，}并非君王的悔改，而是认罪，按照天主的法则的常例，做恶事的人会因自己的宣认而定自己的罪。\BibleQuote{但为了同席的人，}经上说。还有什么比为了避免得罪同席的人而犯下谋杀更卑鄙的呢？接着又说：\BibleQuote{为了他的誓言。}这是何等新奇的宗教！他倒不如发个虚誓。所以主在福音中吩咐我们总不可发誓，\BibleRef{玛 5:34}，以免有发虚誓的缘由，也不需要得罪任何人。因此，一个无辜的人竟被杀死，只为不违背一个誓言。我知道最该憎恶哪一个。暴君的誓言比发虚誓更不可容忍。\par

29. 当有人从宴席跑向监牢时，\BibleRef{谷 6:27}，谁不会以为那是下令释放先知呢？我说，谁听到这是\Person{黑落德}的生日，有盛大的宴席，还给少女任意选择她想要的，会不以为那人是被派去释放\Person{若翰}的呢？残忍与美味有何共通之处？死亡与快乐有何共通？先知在节日时分，被节日的命令匆匆带去受刑，而这一命令本该是用来释放他的；他被剑杀死，头被放在盘子里端上来。这盘子正合他们的残忍，好使他们那无法满足的残暴能饱餐一顿。\par

30. 看哪，最凶残的国王，看看配得上你宴席的景象吧。伸出你的右手，好让你的残忍毫无欠缺，让圣血的溪流从你指间倾泻。既然对如此闻所未闻之残暴的饥渴无法靠宴席满足，也无法靠杯觥解渴，你就饮用从那被砍下头颅仍在流血的血管中涌出的血吧。看哪，那些眼睛，即使在死亡中，仍是你罪行的见证，正转离眼前的佳肴。眼睛在闭合，与其说是因为死亡，不如说是因为对奢华的憎恶。那张失血的金口，其判决你无法忍受，如今沉默了，而你却仍在惧怕。然而那舌头，即使在死后仍惯于如生前般履行职责，虽然颤抖着，仍谴责了乱伦之罪。这颗头被送到\Person{黑落狄雅}那里：她欢喜，她雀跃，仿佛逃脱了罪行，因为她杀死了自己的判官。\par

31. 圣善的妇女们，你们说什么呢？你们可看到你们应教导女儿什么，以及不应教导她们什么吗？她跳舞，但她是淫妇的女儿。但那些端庄、贞洁的妇女，要让她们教导女儿信仰，而非跳舞。而你们，庄重谨慎的男子，要学习避开可憎之人的宴席。若宴席如此，不虔敬者的审判又将是怎样呢？\par

% structure: p0045 source=h2 target=subsection reason=relative-heading-rank; base=h2

\subsection{第7章}

\Emph{为了回复\Person{玛尔切利纳}——她曾问及：对于那些为逃避暴力而自杀的人，该作何想——\Person{圣盎博罗削}叙述了贞女\Person{彼拉奇亚}及其母亲与姐妹的事迹，并接着讲述了他自己的一位祖先、\Person{真福索泰里斯}的殉道。}\par

32. 在我讲话接近尾声时，圣善的姊妹，你提出了一个很好的建议，要我谈谈我们应当如何看待那些人的功劳：他们为免落入迫害者之手，从高处跳下，或投河自尽，因为\WorkTitle{圣经}禁止基督徒自戕。确实，就那些处于必须保守贞洁的处境中的贞女而言，我们有一个明确的答复，因为存在着殉道的实例。\par

33. \Person{圣彼拉奇亚}曾住在\Place{安提约基雅}，年约十五岁，是贞女们的姐妹，自身也是贞女。迫害的消息一传来，她便闭门在家，眼见自己被那些要夺去她信仰与贞洁的人包围，而母亲和姐妹又不在身边，毫无护卫，却愈发充满天主。她对自己说：\Quote{我们该怎么办呢？除非你，你这贞洁的俘虏，多加思量。我既想死，又怕死，因为我并非遇见死亡，而是寻求死亡。若被允许，我们就死吧；即使他们不允许，我们仍然去死。天主并不因对邪恶的补救之策而发怒，信德也允许此举。诚然，若我们思考暴力一词的真义，自愿的行为怎能算暴力呢？反倒想死而不能，那才是暴力。我们不怕任何困难。天下有谁想死而不能，因通往死亡之途如此多而易行？譬如，我现在就可冲向那些亵圣的祭台，将其推翻，用我的血熄灭点燃的火焰。我不怕我的右手无力挥击，也不怕我的胸膛畏缩疼痛。我不会让我的肉身留下罪。我不怕利剑难寻。我能用自己的武器而死，无需刽子手之助，死在我母亲的怀抱里。}\par

34. 据说她装饰了自己的头，并穿上了新娘的礼服，让人以为她是去赴新郎，而非赴死。但当那些可憎的迫害者看到已失却了她贞洁的猎物时，便开始搜寻她的母亲与姐妹。然而她们，以属神的飞升，已占据了贞洁的阵地，当时，一边是迫害者突然逼近，另一边是一条湍急的河流阻断逃路，她们说：\Quote{我们怕什么呢？看这水，有什么能阻止我们领受圣洗呢？这正是那圣洗，借以赦免诸罪，追寻天国。这是圣洗，之后无人再犯罪。让这水接纳我们吧，它惯于使人重生。让这水接纳我们吧，它造就贞女。让这水接纳我们吧，它开启天堂，保护弱者，隐藏死亡，成就殉道者。我们向祢祈求，天主，万有的创造者，莫让这水分散我们失去了生命气息的身体；莫让死亡分离我们的丧礼，我们生前的情感常使我们相联；但愿我们的坚贞如一，我们的死亡如一，我们的埋葬也如一。}\par

35. 说完这些话，她们稍微束起衣襟，既遮掩了端庄又无碍步履，手挽着手如同领舞一般，径直走向河床中央，朝着水流更急、深度更陡之处迈步。无人退缩，无人止步，无人试探落脚之处，她们只在下探着河底时才焦急，在水浅时忧伤，在深水处欢欣。你能看见那虔诚的母亲紧握着双手，为她的寄托而欢悦，生怕跌倒，甚至生怕流水将女儿从她身边冲走。她说：\Quote{基督啊，我将这些牺牲献上，作为贞洁的领袖、我旅途的向导、我苦难的同伴。}\par

37. 然而，谁又会对此感到惊奇：她们生前如此坚贞，死后仍保持着身体姿势不变呢？水没有使她们的遗体裸露，湍急的河流也没有将她们冲走。更有甚者，那位圣善的母亲，虽已毫无知觉，却仍保持着慈爱的紧握，守住她所系的神圣之结，至死未曾松手，为让那位已向信仰还报己债的她，在死时留下自己的虔诚作为嗣业。因为她将那些与自己结合一同殉道的人，甚至认作坟墓中的同伴。\par

38. 但是，我的姐妹，为什么对你使用外族人做例子呢？你是从一位\OriginalTerm{殉道的}{martyred}祖先那里，因血脉相传而得到世代\OriginalTerm{贞洁}{chastity}的感召的。因为你生活在乡间，既无\OriginalTerm{守贞}{virgin}的同伴，又无教导你的老师，你从何处学来的呢？你没有扮演\OriginalTerm{门徒}{disciple}的角色——因为这没有教导是做不到的——你扮演的是德行的继承者。\par

39. 因为圣\Person{索特里斯}——你们家族的祖先——怎么可能不是你的志向的肇始者呢？她在教难的时代，被奴隶们的侮辱激至痛苦的巅峰，竟将她的脸也交给了刽子手——脸在全身受刑时通常是免受伤害的，且更像是观看痛苦而非承受痛苦。她如此勇敢忍耐，以致当她将柔嫩的面颊献给刑罚时，刽子手在打击之前就已失败，而那位\OriginalTerm{殉道者}{martyr}在伤痛下仍未屈服。她不移动她的脸，她不转开她的面容，她不发出一声呻吟或一滴眼泪。最后，当她战胜了其他各种惩罚后，她找到了她所渴望的剑。\par

\SourceRecord{Translated by H. de Romestin, E. de Romestin and H.T.F. Duckworth.；Nicene and Post-Nicene Fathers, Second Series；Vol. 10.；Edited by Philip Schaff and Henry Wace.；Buffalo, NY: Christian Literature Publishing Co.,；1896.；<http://www.newadvent.org/fathers/34073.htm>.}
